% 本文件由 LaTeX 排版 Agent 根据有效 translation.json 自动生成。
% author=John Chrysostom; work_id=2308; title=《弟铎书》讲道集

\chapter{《弟铎书》讲道集}

% structure: p0001 source=h1 target=omitted reason=page-title-already-rendered-by-parent

第一篇讲道\newline{}第二篇讲道\newline{}第三篇讲道\newline{}第四篇讲道\newline{}第五篇讲道\newline{}第六篇讲道\par

\SourceRecord{Translated by Philip Schaff.；Nicene and Post-Nicene Fathers, First Series；Vol. 13.；Edited by Philip Schaff.；Buffalo, NY: Christian Literature Publishing Co.,；1889.；<http://www.newadvent.org/fathers/2308.htm>.}

% structure: p0001 source=h1 target=section reason=page-title

\section{弟铎书讲道（一）}

\BibleRef{铎 1:1-4}\par

\BibleQuote{保禄，天主的仆人，耶稣基督的宗徒，为传播天主的选民的信仰，及认识那合乎虔敬的真理，这真理是在于盼望永生；这永生，那不能撒谎的天主，在万世以前就预许了，但在适当时期藉着宣讲彰显了祂的圣言；这宣讲是照我们救主天主的命令，托付给了我；我现在写信给弟铎，我按共同信仰的真儿子：愿恩宠、怜悯与平安，由天主父及主耶稣基督我们的救主赐予你。}\par

弟铎是保禄的同伴中一位经过考验的人；否则，保禄不会将整个岛的职责托付给他，也不会命令他补充所缺，如他说：\BibleQuote{你要将所缺的补足。} \BibleRef{铎 1:5} 他若对他没有极大的信任，也不会赐予他管辖这么多主教的权柄。据说弟铎当时也是年轻人，因为保禄称他为儿子，尽管这并不能证明。我想在《宗徒大事录》中曾提到他。也许他是格林多人，除非另有同名者。保禄还召唤则纳，并命令将阿颇罗送到他那里，却从未召唤弟铎。\BibleRef{铎 3:13} 因为保禄也为他们的卓越品德和勇气在皇帝面前作证。\par

此后似乎过了一些时间，保禄写这封信时似乎是自由的。因为他没有提到自己的试炼，而是不断强调天主的恩宠，这足以鼓励信徒持守美德。因为知道他们本来应得什么，又被转移到什么状态，且是藉着恩宠，以及所赐予他们的恩惠，这不是小小的鼓励。他也攻击犹太人，如果他责备整个民族，我们不必惊奇，因为他对迦拉达人也同样说：\BibleQuote{糊涂的迦拉达人！} \BibleRef{迦 3:1} 这并非出于吹毛求疵的性情，而是出于爱。如果他是为自己这样做，人们可以公正地责备他；但如果出于他对福音的热忱，那就不算侮辱。基督也多次责备经师和法利塞人，不是为自己，而是因为他们使其他人堕落。\par

他写了一封短信，这合情合理，也是弟铎德行的证明：他不需许多话，只需简短的提醒。但这封信似乎是在写给弟茂德之前写的，因为那封信是他在临终和监禁中所写，而这里他似乎是自由且不受约束的。因为他说的\BibleQuote{我决定在尼苛颇里过冬} \BibleRef{铎 3:12} 证明他当时尚未被捆绑，不像他写给弟茂德时那样。\par

第1节。\BibleQuote{保禄，天主的仆人，耶稣基督的宗徒，为传播天主的选民的信仰。} \BibleRef{铎 1:1}\par

你看他如何不加区别地使用这些词，有时称自己为\Quote{天主的仆人}，有时为\Quote{基督的仆人}，从而表明父与子之间没有区别。\par

\BibleQuote{为传播天主的选民的信仰，及认识那合乎虔敬的真理，这真理是在于盼望永生。}\par

\BibleQuote{为传播天主的选民的信仰。}这是因为你们相信了，还是因为你们被托付了？我认为他的意思是，他被托付了天主的选民，也就是说，不是因我自己的成就，也不是因我的辛劳，我获得了这尊荣。这完全是祂的美善所成就的，祂托付给了我。然而，为使恩宠不显得没有理由（因为仍不完全出于祂，为何不托付给别人？），他因此加上\BibleQuote{及认识那合乎虔敬的真理。}因为我是为这认识而被托付的，或者更确切地说，这也是祂的恩宠所托付给我的，因为祂也是这认识的源头。因此基督自己说：\BibleQuote{不是你们拣选了我，而是我拣选了你们。} \BibleRef{若 15:16} 在别处，这位蒙福者也写道：\BibleQuote{我将认识，如同我被认识一样。} \BibleRef{格前 13:12} 又说：\BibleQuote{倘使我得以达到那为基督耶稣所达到我的目标。} \BibleRef{斐 3:12} 我们先被达到，然后我们认识；我们先被认识，然后我们达到；我们先被召叫，然后我们服从。但他说\BibleQuote{为传播选民的信仰}，这一切都归功于他们，因为我是为他们而做宗徒，不是因我的价值，而是\BibleQuote{为选民的缘故}。正如他在别处说：\BibleQuote{一切都是你们的，无论是保禄，或是阿颇罗。} \BibleRef{格前 3:21}\par

\BibleQuote{及认识那合乎虔敬的真理。}因为在其他事物中也有真理，不是合乎虔敬的；例如农业知识、工艺知识，是真实的真理；但这真理是合乎虔敬的。或者说，这\Quote{根据信仰}，意味着他们相信了，正如其他选民相信并认识真理一样。这认识来自信仰，而非推理。\par

\BibleQuote{这真理是在于盼望永生。}他谈到了现世生命，即在天主恩宠中的生命，也谈到了未来，并在我们面前摆出天主赐予我们的怜悯之后应得的赏报。因为祂愿意加冕我们，因为我们相信了，并从错误中被解救出来。你看，序言充满了天主的怜悯，而整封信尤其如此，以此激励圣人本人及其门徒更加努力。因为没有比常常记念天主的怜悯（无论是公开的还是私下的）更能使我们受益的了。如果我们因接受朋友的恩惠，或听到他们仁慈的话语或行为而心中温暖，那么当我们看到自己陷入何种危险，而天主又怎样拯救我们脱离一切时，我们就更会热心事奉祂。\par

\Quote{并认识真理。}他是就着预像来说这话的。因为那是一个\Quote{认识}和\Quote{虔敬}，却非真理的认识，然而也不是虚谎；它是虔敬，但只是预像和预表。他说得对：\Quote{在永生的希望中。}因为前者是在现世生命的希望中。经上记载：\BibleQuote{遵行这些事的人，必因此而生活。}\BibleRef{罗 10:5}你看他一开始就指出了恩宠的不同。他们不是被选者，而是我们。因为他们若曾被称为被选者，如今却不再这样称呼了。\par

第二节。\BibleQuote{这是那不能撒谎的天主，在万世之前所预许的。}\BibleRef{铎 1:2}\par

这就是说，并非现在改变了心意，而是从起初就这样预定。他常如此断言，如他说：\BibleQuote{被分别出来归于天主的福音。}\BibleRef{罗 1:1}又说：\BibleQuote{祂所预知的人，也预定他们。}\BibleRef{罗 8:29}以此表明我们崇高的根源，即祂不是现在才首先爱我们，而是从起初；从起初就被爱，这并非小事。\par

\Quote{这是那不能撒谎的天主所预许的。}若祂\Quote{不能撒谎}，祂所预许的必然成就。若祂\Quote{不能撒谎}，我们就不该怀疑，即使是在死后。他说：\Quote{这是那不能撒谎的天主，在万世之前所预许的}，借着\Quote{在万世之前}，他也表明这值得相信。这不是因为犹太人没有进来，才预许这些事。这是从起初就计划好的。因此请听他接下来说的：\par

\Quote{但在适当的时候，彰显了。}\par

那么为何延迟呢？出于对人的关怀，并且要在适宜的时候成就。先知说：\BibleQuote{上主，现在是祢行动的时候了}\BibleRef{咏 119:125}。因为\Quote{祂适当的时候}指的是合适、应时、恰当的时候。\par

第三节。\BibleQuote{但在适当的时候，借着宣讲彰显了祂的话；这宣讲是交托给我的。}\BibleRef{铎 1:3}\par

这就是说，宣讲交托给了我。因为这包含了万事：福音、现世与未来之事、生命、虔敬、信德，以及一切。\Quote{借着宣讲}，即公开地、全然放胆地，因为这就是\Quote{宣讲}之意。如同传令官在剧场中当着众人宣告，我们也这样宣讲，不增添什么，只宣告我们所听见的。因为传令官的优点在于向众人宣告真实发生的事，而不增添或删减任何东西。因此，若必须宣讲，就必须放胆直言。否则，就不是宣讲。为此基督没有说\Quote{在屋顶上}告诉，而是说\BibleQuote{在屋顶上宣讲}\BibleRef{玛 10:27}，既以地方也以方式来表明当行之事。\par

\Quote{这是按照我们救主天主的命令交托给我的。}\par

\Quote{交托给我}和\Quote{按照命令}这些说法，表明这事可信，使人不以为可耻，也不迟疑或不满足。既然如此，这是一条命令，就不由我处置。我执行所命令的事。因为当行的事中，有些在我们能力之内，有些则不然。祂所命令的，就不在我们能力之内；祂所允许的，则留给我们选择。例如：\BibleQuote{凡向弟兄说\InnerQuote{傻瓜}的，应受地狱的火。}\BibleRef{玛 5:22}这是一条命令。又说：\BibleQuote{若你献礼物在祭台上，在那里想起你的弟兄有什么怨你的事，就把礼物留在祭台前，先去与弟兄和好，然后来献礼物。}\BibleRef{玛 5:23}这同样是一条命令。但当祂说：\BibleQuote{你若愿意成为成全的，去变卖你所有的}\BibleRef{玛 19:21}，以及：\BibleQuote{能领受的，就领受吧}\BibleRef{玛 19:12}，这就不是命令，因为祂使听者成为这件事的主宰，并让他选择做或不做。因为这些事情我们可以做，也可以不做。但命令并不留给我们选择；我们必须遵守，否则因不遵守而受罚。当他这样说时，也暗示了这一点：\BibleQuote{责任已委托给我；若不传福音，我就有祸了。}\BibleRef{格前 9:16}我将更明确地陈述，使众人明白。例如：受托管理教会、并蒙主教职务尊荣的人，若不向百姓宣告当行之事，就必负责。但平信徒并无此责任。为此保禄也说：\Quote{按照我们救主天主的命令}，我这样做。请看这些称谓如何与我说的一致。因为上面说了\Quote{不能撒谎的天主}，此处则说\Quote{按照我们救主天主的命令}。既然祂是我们的救主，而且祂命令这些事是为了让我们得救，这就不是出于爱发令。这是信德之事，是救主天主的命令。\par

\Quote{致弟铎，我自己的儿子，} 这是指，我的真子。因为人可能不是真子，正如他所说的那人：\Quote{若有人称为弟兄，是淫乱的、贪吝的、拜偶像的、辱骂的、酗酒的，这样的人，不要与他一同进食。}\BibleRef{格前 5:11} 这里有一个儿子，但却不是真子。他确实是儿子，因为他曾接受恩宠，并且重生了；但他不是真子，因为他配不上他的父亲，并且叛逃到另一个篡权者的统辖之下。因为在自然的儿女中，真子与孽子是由生父和生母决定的。但在此事上并非如此，而是取决于心性。因为一个曾是真子的人可能变成孽子，一个孽子也可能变成真子。因为决定因素不是自然的力量，而是选择的权能，因此它常发生变化。敖尼息摩曾是真子，但他又不再是真子，因为他成了\Quote{无益的}；然后他又再成为真子，以致被宗徒称为他的\Quote{心腹}。\BibleRef{费 1:12}\par

\Quote{致弟铎，因共同的信德作我真子。}\BibleRef{铎 1:4}\par

\Quote{因共同的信德}是什么意思？在他称他为自己的儿子，并接受了为父的尊荣之后，请听他如何降低和贬抑那尊荣。他接着说，\Quote{因共同的信德}；这就是说，就信德而言，我并无胜过你的地方；因为它是共同的，你我都因它而重生。那么，他为何称他为儿子呢？要么只是表达对他的爱，要么是他在福音中的优先地位，要么是表明弟铎曾由他光照。为此缘故，他称信友们既是孩子又是弟兄；弟兄，因为他们由同一信德所生；孩子，因为由他亲手所生。因此，他提到共同的信德，暗示他们的弟兄关系。\par

\Quote{愿恩宠与平安，由天主圣父及我们的救主基督耶稣赐予。}\BibleRef{铎 1:4}\par

因为他曾称他为儿子，他接着说，\Quote{由天主圣父}，借此提升他的思想，表明他是谁的儿，并且不仅提到共同的信德，还加上\Quote{我们的父}，暗示他与他有同等的尊荣。也要注意他如何为那教师献上同样的祈祷，如同为门徒和群众一样。事实上，他需要这样的祈祷，与别人一样多，或者更多，因为他要遭遇更大的敌意，并且更容易陷入得罪天主的必要中。因为地位越高，司铎职务的危险就越大。因为在主教职位上，一件善行足以将他提升至天堂，一件过犯足以将他沉入地狱本身。因为，暂且不提日常发生的其他一切情况，倘若他因友谊或其他原因，竟然提升一个不配的人到主教职位，并将一个大城市的治理权交给他，请看他将自己置于何等大火之中！因为他不仅要为那些丧亡的灵魂（因那人的不虔敬而丧亡）交账，还要为那人所做的一切错事负责。因为一个在私人身份中不虔敬的人，一旦被提升到权力地位，就会更加如此。倘若一个虔诚的人在被提升到统治地位后仍能保持虔诚，那真是难能可贵。因为那时他会被虚荣、贪财和任性所猛烈攻击，因为职务赋予他权力；还会被冒犯、侮辱、谴责和无数其他邪恶所攻击。因此，如果任何人不虔敬，他一旦被提升到职位上就会变得更加不虔敬；而任命这样统治者的那个人，要为他所犯的一切罪行和全体人民负责。但如果关于那使人跌倒的，经上记着说：\Quote{若有人使一个小子跌倒，倒不如把驴拴的磨石挂在他的颈上，沉入海底}\BibleRef{玛 18:6}；那么那使这么多灵魂跌倒的，整个城市和居民，众多家庭，男人、女人、儿童、市民和农夫，城市本身和所属一切地方的居民，他将要忍受什么呢？即使说三倍之多，也算不得什么，因为他所面临的报复和惩罚如此严厉。所以主教尤其需要天主的恩宠与平安。因为如果没有这些，他治理人民时，一切都会毁灭和丧失，因为缺乏这些舵。尽管他精通驾驶之术，如果没有这些舵，他也会把船和船上的人都沉没。\Quote{天主的恩宠与平安。}\par

因此，我感到惊讶，那些渴望如此重担的人。可怜而不幸的人，你看到你所渴望的是什么吗？如果你独自一人，默默无闻，即使你犯了一万次错误，你只需为一个灵魂交账，只为他负责。但当你被提升到这一职位时，考虑你将为多少人而受罚。听听保禄所说：\Quote{你们要听从那些领导你们的人，并要服从他们，因为他们为了你们的灵魂常醒寤不寐，好像要交账的人。}\BibleRef{希 13:17} 但你渴望尊荣和权力吗？但这尊荣有什么乐趣呢？我承认，我看不到。因为要做真正的统治者是不可能的，因为这取决于你属下的人是否服从。且对任何细心思考此事的人而言，显然一个主教与其说是来统治，不如说是来服事许多主人，这些主人有相反的心愿和看法。因为一人称赞的，另一人责备；这人指责的，那人赞赏。那么他该听谁，该顺从谁？不可能！那用金钱买来的奴隶，如果主人的命令互相矛盾，他也会抱怨。但如果你在这些多数主人下相反命令时感到悲伤，你甚至为此被定罪，众口一词攻击你。那么告诉我，这是尊荣吗？这是统治吗？这是权力吗？\par

担任主教职的人要求贡献金钱。不愿贡献的人不仅扣留它，而且为了使自己不显得因冷漠而扣留，他就指责他的主教。他说，他是个贼，是个强盗，他吞没了穷人的财物，他吞噬了穷人的权利。停止你的诽谤！你要说这些话到几时？你不愿贡献吗？没有人强迫你，没有强制。你为什么辱骂那劝告和劝导你的人？有人陷入贫困，而那人因无力或其他阻碍没有伸出援手？那人不被谅解，在这种情况下受到的责备比另一种更严厉。那么这就是治理！而他不能报复自己。因为他们是他的内脏，犹如内脏肿胀，尽管给头和身体其余部分带来疼痛，我们也不敢报复，我们不能拿刀刺穿它们；同样，如果我们统治下有人是这种样子，并用这些指控制造麻烦和混乱，我们不敢报复自己，因为这远非为父的心肠，但我们必须忍受痛苦，直到他变得健全和健康。\par

用银钱买来的奴隶有一份指定的工作，完成后他便自主。但主教被各方分心，并被期望做许多超出他能力的事。如果他不会说话，就有很大的怨言；如果他会说话，就被指责为虚荣。如果他不能复活死人，他们就说他毫无价值：他们说，某人虔诚，但这个人不虔诚。如果他吃一顿适中的饭，为此被指责，他们说他该被勒死。如果他在浴室被看见，他受到许多非议。简言之，他不该看见太阳！如果他做同我所做的同样的事，如果他洗澡、吃喝、穿同样的衣服，并管理房屋和仆人，那凭什么他被立在我之上？但他有仆人伺候他，有驴可骑，那为什么他被立在我之上？但你说，难道他应该没有人伺候吗？他应该自己生火、打水、劈柴、去市场吗？这是多大的贬低！就是圣宗徒也不愿任何传道者去照顾寡妇的饭食，他们认为这是不配他们的事：而你竟要让他们降格到你自己的仆人的职务吗？你，命令这些事的人，为什么不来做这些服务？告诉我，难道他为你做的服务不比你的（肉体上的）更好吗？你为什么不派你的仆人去伺候他？基督洗了祂门徒的脚；对你来说，把这项服务给你的老师是件大事吗？但你自己不愿做，还吝啬不给他。难道他该从天上获取生计吗？但天主不愿意这样。\par

但你说：\Quote{宗徒有自由人服侍他们吗？}那么你愿意听宗徒们怎样生活吗？他们走了长途，自由人和尊贵的妇女为他们的帮助献出了生命和灵魂。但听这受祝福的宗徒这样劝勉：\BibleQuote{你们要尊敬这样的人}\BibleRef{斐 2:29}；又说：\BibleQuote{因为他为了基督的工作，至于死亡，不顾自己的性命，为要弥补你们对我的服事之缺欠。}看他所说的话！但你对你的神父一句话也舍不得，更不用说为他冒任何危险了。但你说：\Quote{他不该常去浴室。}这在哪里被禁止？不洁净并无可夸之处。\par

这些根本不是我们所指责或称赞的事。因为主教所应具备的品性是不同的，如无可指摘、节制、端庄、好客、善于教导。宗徒要求这些，我们也应当在教会的统治者身上寻找这些，而非别的。你并不比保禄更严格，或者不如说比圣神更严格。如果他好斗、粗暴、残忍、不慈悲，就指责他。这些事与主教不相配。如果他奢侈，这也是可责备的。但如果他照顾自己的身体以便为你服务，如果他注意健康以便有用，他因此就该被指责吗？你不知道身体上的软弱与灵魂上的软弱一样，既损害我们自己也损害教会吗？否则，保禄为何在处理此事时，写信给弟茂德说：\BibleQuote{为了你的胃和你屡次的疾病，用一点酒吧}\BibleRef{弟前 5:23}？因为如果我们能仅凭灵魂修德，我们就不必照顾身体。那么我们为何出生呢？但如果身体贡献了很大一部分，忽视它岂不是极端的愚昧吗？\par

设想一个人被授以主教职，并受托担任教会的公共职务，他在其他方面有德性，并拥有一切司铎应具备的品性，但若因重病而常年卧床，他能提供什么服务？他能执行什么使命？他能进行什么探访？他能责备或劝诫谁？我说这些话，是要你们学会不无故指责他，反而善意接待他；也是为了让任何渴望在教会中掌权的人，看到随之而来的种种辱骂，从而熄灭那欲望。这样的职位实在危险很大，它需要\Quote{天主的恩宠与平安}。为使这恩宠与平安丰盛地赐予我们，请为我们祈祷，我们也为你们祈祷，好使我们正确修德，获得所应许的福分，借着耶稣基督，偕同祂，等等。\par

\SourceRecord{Translated by Philip Schaff.；Nicene and Post-Nicene Fathers, First Series；Vol. 13.；Edited by Philip Schaff.；Buffalo, NY: Christian Literature Publishing Co.,；1889.；<http://www.newadvent.org/fathers/23081.htm>.}

% structure: p0001 source=h1 target=section reason=page-title

\section{\WorkTitle{提铎书讲道 第二篇}}

铎 1:5,6\par

\BibleQuote{我为此把你留在克里特，是要你将那些欠缺的事安排妥当，并照我所吩咐的，在各城设立长老：若有人无可指摘，只作一个妻子的丈夫，有信德的儿女，不被控告放荡或不守规矩。}\par

古代人的生活是行动与奋斗；相反，我们的生活却是懒散。他们知道，人被带到世上是为了这个目的：要按那带他们入世者的旨意而劳苦；但我们，仿佛被安置在这里只是为了吃、喝，过享乐的生活，不关心属灵的事。我不仅指宗徒们，也指那些跟随他们的人。你看他们因此走遍各地，将此作为他们唯一的事业，完全如同在异乡生活，如同在地上没有常城的人。因此，请听蒙福的宗徒所说的话：\par

\BibleQuote{我为此把你留在克里特。}\par

仿佛整个宇宙是一座房屋，他们将它分给自己，在各地管理其事务，每人照管自己的那一份。\par

\BibleQuote{我为此把你留在克里特，是要你将那些欠缺的事（修订版作\Quote{是欠缺的}）安排妥当。}\par

他并不是以命令的方式吩咐这事，而是说\BibleQuote{你要安排妥当}。这里我们看到一个灵魂完全摆脱嫉妒，到处寻求他门徒的利益，并不好奇地关切这好处是由自己还是由别人做成的。因为凡有危险和极大困难之处，他亲自处理；而那些较为荣耀和值得称赞的事，他则委托给他的门徒，如主教的任命，以及其它需要进一步安排，或可说需要更完善的事。你说什么？难道他进一步安排你的工作，而你不认为这是使你觉得羞愧的耻辱吗？绝非如此；因为我只着眼于公共利益，不论是由我做成还是由别人做成，对我都一样。凡在教会中做主的人，应当有这样的心境：不是寻求自己的荣耀，而是公共利益。\par

\BibleQuote{在各城设立长老}，这里他是指主教，如同我们以前说的，\BibleQuote{照我所吩咐的}。\BibleQuote{若有人无可指摘}。\BibleQuote{在各城}，他说，因为他不想整个岛托付给一个人，而是每人有自己的职责和照顾，这样他自己也少些劳苦，他治下的人也得到更多关注，因为教师不必奔走管理许多教会，而是只专注于一个，将它整顿好。\par

\BibleQuote{若有人无可指摘，只作一个妻子的丈夫，有信德的儿女，不被控告放荡或不守规矩。}\BibleRef{铎 1:6}\par

他为什么提出这样的人？为了堵住那些谴责婚姻的异端者的口，表明婚姻本身并非不圣洁，而是如此可敬，以致一个已婚之人可以登上圣座；同时责备那些放荡之人，不允许那些再婚的人进入这一高位。因为一个对已故妻子毫无敬意的人，怎能成为一个好的管理者？他又会招致什么谴责呢？因为你们都知道，虽然法律并未禁止第二次婚姻，但这件事情容易引起许多非议。因此，他希望一个统治者不让受他统治的人有指责的把柄，所以他这样说：\BibleQuote{若有人无可指摘}，也就是说，若他的生活无可指责，若他没有给任何人中伤他品行的机会。请听基督的话：\BibleQuote{如果你内的光明成了黑暗，那黑暗是多么大啊！}\BibleRef{玛 6:23}\par

\BibleQuote{有信德的儿女，不被控告放荡或不守规矩。}\par

我们应该注意他对儿女的关心。因为连自己的儿女都不能教导的人，怎能教导别人呢？如果他不能管束那些从起初就与他在一起、由他抚养、并因法律和本性他有权管理的人，他又怎能有益于外人呢？因为如果父亲的失职不是很大，他决不会让那些从一开始就在他权下的人变坏。因为不可能，确实不可能，一个受到极大照顾、得到很多关注的人会变坏。罪恶并不如此占优势，以致能战胜如此多的预先照顾。但如果他忙于追求财富，把儿女当作次要的事，没有给予他们很多关心，那么他同样不配。因为如果当本性驱使他时，他竟如此没有感情或如此麻木，认为财富比儿女更重要，那么他又怎能被提升到主教之位和如此大的管理权呢？因为如果他不能约束他们，这是他软弱的巨大证明；如果他不关心，他的缺乏感情就大受责备。那么，忽视自己儿女的人，又怎能照顾别人家的儿女呢？他不仅说\BibleQuote{不放荡}，而且甚至\BibleQuote{不被控告放荡}。不能有恶名，或这样的看法。\par

\BibleQuote{主教必须是无可指摘的，如同天主的管家；不任性，不轻易发怒，不酗酒，不斗殴。}\BibleRef{铎 1:7}\par

因为一个外邦的统治者，由于他依据法律和强制来统治，或许不征求被统治者的意见。但那应当以人民自愿顺从的方式统治的人，如果他的行为事事按自己的意愿，不与任何人商议，那么他的治理就变成暴君式的，而非民主的。因为他必须是\BibleQuote{无可指摘的，如同天主的管家；不任性，不轻易发怒}。因为一个连自己都未曾教导好如何驾驭怒气的人，怎能教导别人呢？因为权力引来许多诱惑，它使人更加严厉、难以取悦，即使原本很温和的人，周围也环绕着如此多的触怒之事。如果他事先没有在这美德上操练自己，他就会变得严厉，并伤害和毁灭很多他管辖下的事物。\par

\Quote{不可嗜酒，不可擅斗。}他在这里说的是狂妄的人。因为他当以劝告或斥责行事，而非蛮横。何必，请告诉我，要侮辱呢？他应当恐吓，应当警醒，以地狱的威胁深入灵魂。但受侮辱的人变得更无耻，反而轻视侮辱他的人。没有什么比侮辱更招致轻视；它使侮辱者蒙羞，阻止他得到应有的尊重。他们的言论当十分谨慎。在指责罪过时，他们应记起未来的审判，但避免一切蛮横。然而，若有人阻止他们尽职，他们必须用全权追究此事。\Quote{不可擅斗，}他说。教师是灵魂的医生。但医生不打击，反医治和恢复那打击他的人。\Quote{不可贪求不义之财。}\par

第8节。\BibleQuote{却要爱好款待客旅，爱好行善，端庄，公义，圣善，节制。}\BibleRef{铎 1:8}\par

第9节。\BibleQuote{持守那按所受教导的忠信的话。}\BibleRef{铎 1:9}\par

你看他要求多么高的德行。\Quote{不贪不义之财}即表示对金钱极其轻视。\Quote{款待客旅，爱好行善，端庄，公义，圣善}；他意指将一切财物分施给有需要的人。\Quote{节制}；他这里不是指那守斋的人，而是指那能控制自己情欲、口舌、双手、双眼的人。因为这就是节制：不被任何情欲所牵引。\par

\Quote{持守那按所受教导的忠信的话。}此处\Quote{忠信}意指\Quote{真实的}，或那藉着信德所传的，不需要推理或争辩。\par

\Quote{持守}，即关怀它，以此为业。那么，若他不懂外界的学问，怎么办？为此，他说：\Quote{忠信的话，按教导。}\par

\Quote{使他既能劝勉，又能驳倒那些反对的人。}\par

所以需要的不是言辞的浮华，而是坚强的头脑、对圣经的熟练和有力的思想。你没有看见\Person{保禄}征服了全世界，比\Person{柏拉图}和所有其余的人更有力吗？但你是说，那是靠奇迹。不，不单靠奇迹，因为如果你细读\WorkTitle{宗徒大事录}，你会发现他在行奇迹之前也往往凭教导取胜。\par

\Quote{使他能以健全的道理劝勉}，即留住自己的人，并打倒敌人。\Quote{驳倒那些反对的人。}因为若不如此，一切都失丧了。那不知道如何与敌人争战，如何\Quote{将一切思想俘虏，使之顺从基督}，如何击倒推理的人，那不知道就正确道理当教导什么的人，教师的宝座远离他。因为其他品德可在他的属下中找到，例如\Quote{无可指摘，使子女服从，好客，公义，圣善。}但教师的特点在于能教导言语，而现今无人关注这一点。\par

第10节。\BibleQuote{因为有许多人不服约束，空谈欺骗者，尤其是那些受割礼的人；}\BibleRef{铎 1:10}\par

第11节。\BibleQuote{这些人的口必须堵住。}\BibleRef{铎 1:11}\par

你看见他是如何显出他们是这样的吗？因为他们不愿受管辖，却愿管辖人。他暗示了这一点。所以当你不能说服他们时，不必委派他们职务，而要堵住他们的口，为有益于他人。但这有什么好处，如果他们将不服从或不服约束？那么他为何要堵住他们的口呢？为了使其他人从中受益。\par

\Quote{他们为了不义之财，教导不当教的事，颠覆整个家庭。}\par

因为如果一个人承担了教师的职务，却不能与这些敌人作战，不能堵住那些如此无耻之人的口，那么在每个灭亡者的丧亡上，他都成了原因。若有人曾这样劝告：\BibleQuote{不要寻求作审判官，除非你能消除不义}\BibleRef{德 7:6}；我们在此更可以说：\Quote{不要寻求作教师，若你与这职务的尊严不相称；但即使被勉强，也要拒绝。}你看见吗？爱权、贪不义之财是这些恶行的原因。\Quote{他们为了不义之财，教导不当教的事，}他说。\par

伦理。因为没有什么不被这些情欲所败坏。但当暴风落在平静的海上，将海从深处翻起，使泥沙与波浪相混时，同样，这些情欲袭击灵魂，颠倒一切，使心灵的视线昏暗，尤其是贪图虚荣的疯狂。对于钱财的轻视，任何人可能容易达到；但要轻视来自众人的荣誉，则需要巨大的努力，一种哲学家的心境，某种达到至高天堂的属神灵魂。因为没有哪一种情欲如此专横，如此普遍，虽程度不同，但无处不在。那么我们如何制伏它呢？虽不能完全，但至少可以稍微？藉着仰视天空，将天主置于眼前，思想超越尘世的事物。当你渴望荣誉时，设想你已经得到它，并注视其结局，你会发现它毫无价值。思考它所带来的损失：它要剥夺你多少何等大的恩惠。因为你将承受辛劳和危险，却失去它们的果实和酬报。思考大多数人都是恶的，轻看他们的意见。对每一个人，思考他是怎样的人，你就会看到荣誉是多么可笑的事，它更应被称为耻辱。\par

此后，请将你的思想提升到天上的剧场。当你行善时，若认为应当在人前显示，寻求观看者，并费尽心思要被看见，就要思想天主注视着你，那一切愿望便会熄灭。离开地上，注视那在天上的剧场。如果人们称赞你，但日后他们会责备你、嫉妒你、攻击你的名誉；或者即使不如此，他们的称赞于你也无益。天主却不是这样。祂乐于称赞我们的善行。你说得好，博得了喝彩？你得到了什么？如果那些喝彩的人获益、心意更新、成为更好的人、离弃恶行，那么你实在可以欢喜，不是因所得的称赞，而是因这奇妙的转变。但如果他们继续称赞和大声喝彩，却从所称赞的事上未得任何益处，你倒应当忧愁，因为这些事将归于他们的审判和定罪。你却因你的虔诚而得荣耀。若你真虔诚，且自觉无罪，你应当欢喜，不是因被视为虔诚，而是因你确是虔诚。但若你并非如此，却渴望群众的赞誉，要想到在最后的日子审判你的不是他们，而是那完全知晓隐藏之事的天主。而且，如果你虽自觉有罪，却被众人视为纯洁，你非但不应欢喜，反要悲痛哀伤，常常想着那一天，那时一切都要显露，黑暗的隐密事将被揭示出来。\par

你享荣誉吗？拒绝它，因你知道它使你成为负债者。无人尊敬你吗？你应当为此欢喜，因为天主不会将这列入你的账目：你曾享受荣誉。你难道没看见，天主通过先知责备以色列子民时，也提到了这一点：\Quote{我从你们儿子中选立先知，从你们青年中选立献身者}？\BibleRef{亚 2:11}因此你至少会获得这一益处：你不会加重你的惩罚。因为在今生不受尊敬、被人轻视、不受重视、反遭侮辱和蔑视的人，至少获得这一益处（即便没有别的）：他不必因受同仆的尊敬而交账。在其它许多方面他也因此获益：他被降卑和谦逊，即使他愿意，也不能心高气傲，因为他更注意自己。而那享受更多荣誉的人，除了负担巨大的债务外，还会被举扬到傲慢和虚荣之中，成为人的奴隶；随着这种暴政的增强，他被迫做出许多他不愿做的事。\par

因此，知道没有荣耀比拥有荣耀更好，我们不要寻求荣誉，反而要在荣誉被给予时避开它，我们要抛弃它，熄灭那欲望。这话我们同时是对教会管理者和被管理者说的。因为一个渴望荣誉和受尊崇的灵魂不能看见天国。这不是我自己的话。我说的不是自己的话，而是天主圣神的话。他不能看见它，即使他修行美德。因为祂说：\Quote{他们已经获得了他们的赏报。}\BibleRef{玛 6:5}那么，那没有赏报可领的人，如何能看见天国呢？我不禁止你渴望荣耀，但我希望那是真正的荣耀，即来自天主的荣耀。经上说：\Quote{他的称赞不是来自人，而是来自天主。}\BibleRef{罗 2:29}让我们在暗中虔诚，不被排场、炫耀和伪善所累。让我们脱去羊皮，而成为羊本身。没有什么比人的荣耀更无价值。你若看见一群小孩子，还在吃奶的婴孩，你会渴望从他们那里得荣耀吗？对于人的荣耀，你当以同样的心态对待一切人。\par

正因如此，它被称为\OriginalTerm{虚荣}{vainglory}。你看见舞台演员戴的面具吗？何等美丽、灿烂，被塑造得极其优雅。你能给我看任何真实的这样面容吗？绝不能。那么，你曾爱上这些面具吗？没有。为什么？因为它们是空虚的，模仿美丽，而不是真正的美丽。同样，人的荣耀是空虚的，是对荣耀的模仿，不是真正的荣耀。唯独那内在的自然美丽是持久的；那外在加上的往往掩盖了丑陋，只掩盖到晚上日暮。但当剧场解散，面具除去，每个人才显出他的本来面目。\par

因此让我们追求真理，不要像在舞台上演戏一样。告诉我，被众人注目有什么益处？那是\OriginalTerm{虚荣}{vainglory}，别无其他。因为你回到家，独处时，一切立即消失。你去了市场，吸引了所有在场者的目光。你得到了什么？什么也没有。它消散了，像烟一样消散。难道我们爱这样不实在的东西吗？这是多么不合理！何等疯狂！我们只看一件事：永远不寻求人的称赞；但如果它临到我们，我们要轻视、嘲笑、拒绝它。我们要像那些渴望金子却得到泥土的人一样。不要让任何人称赞你，因为它无益；而若有人责备你，也无害。但在天主那里，称赞和责备都带来真实的得失，而来自人的一切皆是虚空。在此我们相似天主：祂不需要从人得荣耀。基督说：\Quote{我不受从人来的光荣。}\BibleRef{若 5:41}告诉我，这是小事吗？当你不愿轻视荣耀时，就说：\Quote{因轻视它，我将相似天主}，你便会立即轻视它。但荣耀的奴隶不可能不是万物的奴隶，比真正的奴隶更为奴役。因为我们并不给我们的奴隶加派这样的任务，如同荣耀向她的俘虏所加派的。她使他们所说、所做、所忍受的都是卑鄙可耻的事，而当她看见他们顺从时，便更急迫地命令。\par

让我们逃走，我恳求你们，让我们从这奴役中逃走。但我们怎能做到呢？如果我们认真思考这世上的事物，如果我们观察到眼前的一切不过是梦、影，没有什么更好的东西；我们就能轻易克服这种欲望，无论是在小事还是大事上都不会被它俘虏。但如果我们在小事上不轻视它，在最重要的事上我们就很容易被它战胜。因此，让我们远离它的根源，那就是愚昧和心胸狭窄，这样，如果我们怀有崇高的精神，就能超越群众的尊敬，将目光投向天堂，并获得那里的美善之物。愿天主恩赐我们众人都有分于此，藉着我们的主耶稣基督的恩宠和慈爱，与祂一起，等等。\par

\SourceRecord{Translated by Philip Schaff.；Nicene and Post-Nicene Fathers, First Series；Vol. 13.；Edited by Philip Schaff.；Buffalo, NY: Christian Literature Publishing Co.,；1889.；<http://www.newadvent.org/fathers/23082.htm>.}

% structure: p0001 source=h1 target=section reason=page-title

\section{论《弟铎书》第三篇讲道}

\BibleRef{铎 1:12-14}\par

\Quote{他们中有一个人，本地的先知，说：\InnerQuote{克里特人常说谎，是恶兽，是懒惰的肚腹。}这个见证是真实的。因此你要严厉地责备他们，使他们在信德上纯全无疵，不听犹太人荒谬的传说和那转离真理之人的诫命。}\par

这里有几个问题。第一，说这话的是谁？第二，保禄为何引用它？第三，他为何提出一个不正确的见证？让我们先作一些预备，再适时解答这些问题。因为当保禄向雅典人讲话时，在演说中他引用了这些词句：\Quote{献给未识之神}；以及：\Quote{我们也是祂的子孙，就如你们自己的诗人所说的。} \BibleRef{宗 17:23} 说这话的是厄庇墨尼德斯，他本人是克里特人，有必要提到他为何受感动说这话。事情是这样的：克里特人有一座朱庇特的坟墓，上面刻着铭文：\Quote{此地安葬着赞，他们称之为朱庇特。} 因此，诗人为了嘲笑克里特人说谎，在行文中又加入这一段以增加讽刺。\par

\Quote{因为他们为你造了一座坟墓，啊，君王，\newline{}你从未死去，却永远存在。}\par

如果这个见证是真的，看哪，多么大的困难！因为如果诗人说他们撒谎是对的，而他们在断言朱庇特会死这件事上撒谎，那么正如宗徒所说的，这是一件可怕的事！亲爱的，请十分精确地注意。诗人说克里特人撒谎，因为他们说朱庇特死了。宗徒证实了他的见证：因此，按照宗徒，朱庇特是不朽的；因为他说：\Quote{这个见证是真实的}！那么，我们该说什么呢？或者更确切地说，我们该如何解决这个问题？宗徒并没有说这个，他只是简单直接地应用这个见证来说明他们惯于说谎。否则他为什么不加上：\Quote{因为他们为你造了一座坟墓，啊，君王}？所以宗徒并没有说这个，而只是说有人说得对：\Quote{克里特人常说谎。} 但我们不仅从这里确信朱庇特不是神，还能从许多其他论据证明这一点，而不是从克里特人的见证。此外，他并没有说在这件事上他们是撒谎的。而且，很可能他们在这个观点上也受了骗，因为他们相信其他神，因此宗徒称他们为说谎者。\par

至于问题：他为何引用希腊人的见证？这是因为当我们从他们自己的作者那里提出见证和指责时，最能使他们困惑，当我们使他们自己所崇拜的人成为指控者时。为此，他在别处引用了\Quote{献给未识之神}。因为雅典人并非从起初就接受所有他们的神，而是不时接纳一些其他的，如来自许珀耳玻瑞亚人的神、对潘的崇拜，以及大小奥秘，所以这些人猜想除了这些之外可能还有某位他们不认识的神，为了虔敬，也为祂立了一座祭坛，刻着\Quote{献给未识之神}，几乎暗示\Quote{如果有一位他们不认识的某位神}。因此他对他们说：你们预先承认的那位，我向你们宣告。至于\Quote{我们也是祂的子孙}，是引用了阿拉托斯的诗句，他之前说：\Quote{大地之路充满朱庇特，大海充满}——然后又加了：\Quote{因为我们也是祂的子孙}，我认为这表明我们来自神。那么保禄如何将关于朱庇特的话转用于宇宙的天主？他并没有把属于朱庇特的转移给天主。而是那些适用于天主、却既不合理也不正当地应用于朱庇特的，他归还给天主，因为天主的名字只属于祂自己，不能正当地给予偶像。\par

他该从什么作者向他们讲话？从先知？他们不会相信。因为面对犹太人，他也不是从福音书论证，而是从先知书。为此他说：\Quote{向犹太人，我就成为犹太人；向没有法律的人，我就成为没有法律的人；向在法律之下的人，我就成为在法律之下的人。} \BibleRef{格前 9:20} 同样，天主也是如此，例如在贤士的事上，祂没有用天使、先知、宗徒或传福音者引领他们，而是用什么？用一颗星。因为他们的艺术使他们熟悉这些，祂就用这样的方式来引导他们。又如，在拉约柜的牛的事上：\Quote{若是上行的路回到自己的境界，那么祂就给我们降这大灾} \BibleRef{撒上 6:9}，正如他们的先知所建议的。那么这些先知说真话吗？不；但祂用他们自己的口驳斥和困惑他们。又如，在巫婆的事上，因为撒乌耳信她，祂就让她宣告将要临到他身上的事。那么为什么保禄堵住了那个鬼神的嘴，它说：\Quote{这些人是至高天主的仆人，对你们传讲救恩的道路} \BibleRef{宗 16:17}？为什么基督禁止魔鬼谈论祂？这是有理由的，因为神迹正在发生。在这里，不是一颗星宣告祂，而是祂自己；而且鬼也不是被崇拜的，因为不是像在说话，所以不应禁止。祂也允许巴郎祝福，没有阻止他。因此，祂处处俯就。\par

这有什么可惊奇的呢？因为祂容许错误的、与祂不相称的见解流行，如祂从前是形体、是可见的。为反对这种见解，祂说：\BibleQuote{天主是神}\BibleRef{若 4:24} 又说，祂喜悦祭献，这远非祂的本性。祂说出与自己宣告相矛盾的话，还有许多类似之事。因为祂从不考虑自己的尊荣，而只考虑什么对我们有益。如果一个父亲不顾自己的尊严，用牙牙学语与孩子们说话，不叫他们的食物和饮料希腊名，而用某些幼稚、粗野的词语，何况天主呢？就连在责备时，祂也俯就，如同祂藉先知所说：\BibleQuote{一个民族岂能更换他们的神？}\BibleRef{耶 2:11}，在整部圣经中，有无数祂在言语和行动上俯就的例子。\par

第13节。\BibleQuote{为此，你要严厉地责备他们，使他们在信德上健全。}\BibleRef{铎 1:13}\par

他说这话，是因为他们的性情乖僻、欺诈、放荡。他们有这无数恶习；又因为他们易于撒谎、欺骗、贪吃、懒惰，所以需要严厉的责备。因为这样的性格不能用温和来管理，\Quote{因此你要责备他们。} 他这里不是论及外邦人，而是论及自己的子民。\Quote{严厉地。} 给他们一个切深的打击。因为不能用一种方法对待所有人，而应按照他们不同的性格和性情分别对待。他这里没有诉诸劝告。因为正如以严厉对待温顺诚实的人可能毁灭他们，同样，奉承一个需要严厉的人，会使其灭亡，不让他被挽回。\par

\Quote{使他们在信德上健全。}\par

这就是健全：不引入任何虚伪或异类的东西。但如果连那些对于食物小心翼翼的人也不健全，而是病弱；因为他说：\BibleQuote{信德软弱的，你们要接纳，但不要为疑惑辩论。}\BibleRef{罗 14:1} 那么，对于那些守同样的斋戒（与犹太人一起）、守安息日、常去他们视为圣的地方的人，又该怎样说呢？我指的是在达佛涅那个地方，被称为玛特罗纳洞穴的那个，以及基里基雅那片被称为撒杜恩平原的地方。这些人如何是健全的呢？他们需要更重的打击。那么，为何他对罗马人不这样做呢？因为他们的性情不同，他们品格更高贵。\par

第14节。他说：\BibleQuote{不要听从犹太人的荒诞传说。}\BibleRef{铎 1:14}\par

犹太人的教条是荒诞传说，有两个原因：因为它们是模仿，又因为时机已过；因为这类事物最终成了荒诞传说。当一件事不该做，做了反而有害时，它就变成荒诞传说，因为无用。所以，如同那些不该被重视，这些也不该被重视。因为这不是健全。如果你相信信德，为何还要添加别的东西，好像信德不足以使人成义？为何因受法律的束缚而奴役自己？难道你对自己所信的不信任吗？这是不健全、不信的心的标志。因为一个忠实的人不怀疑，但这样的人显然怀疑。\par

第15节。他说：\BibleQuote{在洁净的人，一切物都是洁净的。}\BibleRef{铎 1:15}\par

你看，这话是为特定目的而说的。\par

\Quote{但在污秽不信的人，什么也不洁净。}\par

那么，事物洁净或不洁净并非因它们自身的本性，而是因取用者的心态。\par

\Quote{就连他们的理智和良心也都玷污了。}\par

第16节。\BibleQuote{他们声称认识天主，但在行为上却否认祂，是可憎恶的、悖逆的，在一切善事上可废弃的。}\BibleRef{铎 1:16}\par

那么猪是洁净的。为何它被禁止作为不洁呢？它并非由本性不洁；因为\Quote{一切物都是洁净的}。没有比鱼更不洁的，因为它甚至以人肉为食。但它却被允许，视为洁净。没有比鸟更不洁的，因为它吃虫；也没有比鹿更不洁的，据说它的名字来源于吃蛇。然而这些都被吃了。那么为何猪被禁止，以及许多其他东西呢？不是因为它们不洁，而是为了节制过度的奢侈。但若这样说，他们不会信服；因此他们被不洁的恐惧所约束。因为请告诉我，如果我们仔细查考这些事，有什么比酒更不洁？或比水更不洁？他们多用它们来净化自己。他们不摸死人，却用死物来洁净自己，因为祭牲是死的，他们用它来洁净。所以，这是给孩童的教训。在酒的构成中，难道没有粪便的成分吗？因为葡萄树从地中吸收水分，也从投在地上的粪肥中吸收。简而言之，如果我们想非常挑剔，则一切都不洁；如果我们不想挑剔，则没有不洁。然而一切物都是洁净的。天主没有造任何不洁之物，因为除了罪以外，没有不洁的。因为罪达于灵魂，玷污它。其他的不洁是人的偏见。\par

\Quote{但在污秽不信的人，什么也不洁净；就连他们的理智和良心也都玷污了。}\par

因为在洁净的人中，怎么会有不洁之物呢？但灵魂软弱的人使一切变不洁，若对洁净与否进行挑剔查考，他将什么也不触摸。因为按照他们的观念，这些事甚至也不洁净（我指的是鱼和其他东西）；（因为他说：他们的理智和良心都玷污了）但一切都不洁。然而保禄不这样说；他将整个问题转向他们自身。因为他说，没有不洁的，只有他们自己，他们的理智和良心，才是不洁的；没有什么比这些更不洁；但邪恶的意志是不洁的。\par

\Quote{他们声称认识天主，但在行为上却否认祂，是可憎恶的、悖逆的，在一切善事上可废弃的。}\par

第2章第1节。\BibleQuote{至于你，你要讲那合乎健全道理的事。}\BibleRef{铎 2:1}\par

这就是不洁。他们自身是不洁的。但你不可因此缄默。尽你的本分，尽管他们可能不接受你。劝告并忠告他们，尽管他们可能不被说服。这里他更严厉地责备他们。因为那些疯狂的人以为没有什么是静止的，但这并非来自所见的对象，而是来自观看的眼睛。因为他们不稳定且眩晕，他们以为大地随着他们旋转，而大地并不旋转，而是稳固不动。错乱在于他们自身的状态，并非来自元素的任何影响。同样，当灵魂不洁时，它以为万物都是不洁的。因此，拘泥于仪节并非纯洁的标志，纯洁的标志是在一切事上大胆。因为本性纯洁的人敢于面对一切，而玷污的人则不敢面对任何事。我们可以用这一点来反驳马西昂。你看见了吗？超越一切玷污是纯洁的标志，什么也不碰则意味着不洁。这甚至适用于天主。祂取了肉身是纯洁的证明；如果祂因恐惧而未取肉身，那就会有玷污。不吃看似不洁之物的人，自身是不洁和软弱的；吃的人则不是。我们不要称这样的人为纯洁的，他们是不洁的。纯洁的人，是敢于吃用一切食物的人。我们应当对玷污灵魂的事物施加这种谨慎。因为那才是不洁，那才是玷污。这些事物中没有一样是如此。那些味觉败坏的人以为摆在面前的食物是不洁的，但这是他们失调的结果。因此，我们应当明了洁净之物与不洁之物的性质。\par

道德。那么，什么是不洁？罪、恶毒、贪婪、邪恶。正如经上所载：\BibleQuote{你们要洗涤，自洁，从我眼前除掉你们的恶行。}\BibleRef{依 1:16} \BibleQuote{天主，求祢在我内里造一颗纯洁的心。}\BibleRef{咏 51:10} \BibleQuote{离开，离开，从那里出去，不要触摸不洁之物。}\BibleRef{依 52:11} 这些仪节是净化的象征。\BibleQuote{不要触摸死尸，}经上说。因为罪正是如此，是死亡且令人厌恶的。\BibleQuote{癞病人是不洁的。}因为罪是一种癞病，多样而多形。而这些具有这种含义，从下文可以看出。因为如果癞病是全面性的，蔓延全身，他就是洁净的；如果是局部的，他就是不洁的。因此你看见那多变而可变的，就是不洁的。再有，遗精的人是不洁的，这暗指那在灵魂中丢弃种子的人。未受割损的人是不洁的。这些东西不是寓意的，而是预象的，因为那不割除心中邪恶的人，就是不洁的。在安息日工作的人应被处死，这就是说，那并非时时献身于天主的人，必要灭亡。你看见不洁有多少种类。产妇是不洁的。然而天主创造了生育和交合的种子。那么为什么产妇是不洁的，除非另有暗示？那是什么呢？祂意在灵魂中培养虔敬，并阻止其淫乱。因为如果生了孩子的人是不洁的，那么犯了淫乱的人更是不洁的。如果亲近自己的妻子并非全然纯洁，那么与别人的妻子交合更是如此。参加葬礼的人是不洁的，那参与战争和杀戮的人更是如此。如果需要一一列举，会发现许多种不洁。但这些现在不再要求我们。一切都转移到了灵魂。\par

因为肉体的事物离我们更近，所以他由此引入教训。但现在不是这样了。因为我们不应局限于预象和阴影，而应持守真理并维护它：罪是不洁之物。让我们逃避它，让我们戒绝它。\BibleQuote{你走近它，它必咬你。}\BibleRef{德 21:2} 没有什么比贪婪更不洁。这从何显明？从事实本身。因为它玷污了什么？双手、灵魂，甚至那存放不义之财的房屋。但犹太人对此不以为然。然而梅瑟带走了若瑟的骸骨。三松喝了驴腮骨流出的水，并从狮子身上取了蜜，厄里亚由乌鸦和一个寡妇供养。请告诉我，如果我们对这些事吹毛求疵，还有什么比我们的书更不洁？它们是动物皮制成的。那么，淫乱者并非唯一不洁的，还有比他更甚的，如通奸者。但两者都是不洁的，不是由于交合本身（因为按照那个逻辑，与自己的妻子同居的人也会是不洁的），而是由于行为的邪恶以及对邻人最切身利益的损害。你看见了吗？邪恶才是不洁的。有两个妻子的人并不不洁，有众多妻子的达味也不不洁。但当他非法占有一个时，他变得不洁了。为什么？因为他损害并欺骗了邻人。淫乱者不是因交合而不洁，而是因交合的方式，因为它伤害了那女人，并且他们相互伤害，使女人成为公共的，颠覆了自然规律。因为女人应当是一个男人的妻子，因为经上说：\BibleQuote{天主造了他们一男一女。}\BibleRef{创 1:27} 又说：\BibleQuote{二人成为一体。}不是\BibleQuote{那许多人，}而是\BibleQuote{二人成为一体。}这里便有不义，因此那行为是邪恶的。再者，当愤怒超过适当限度时，它使人不洁，不是愤怒本身，而是因它的过度。因为经上不是说：\BibleQuote{那发怒的，}仅仅，而是\BibleQuote{无故发怒的。}因此，每件事上过度贪求是不洁的，因为它出自贪婪和非理性的禀性。因此，我恳求你们，让我们保持清醒，让我们保持纯洁，那真正的纯洁，好使我们堪当看见天主，借着我们的主耶稣基督，偕同祂，永世无穷。\par

\SourceRecord{Translated by Philip Schaff.；Nicene and Post-Nicene Fathers, First Series；Vol. 13.；Edited by Philip Schaff.；Buffalo, NY: Christian Literature Publishing Co.,；1889.；<http://www.newadvent.org/fathers/23083.htm>.}

% structure: p0001 source=h1 target=section reason=page-title

\section{论弟铎书第四讲}

\BibleRef{铎 2:2-5}\par

\BibleQuote{劝老年人要有节制、庄重、审慎，在信德、爱德、忍耐上纯全无缺。同样，劝老年妇人在行为上要合乎圣洁，不诽谤人，不沉湎于酒，是善事的导师，好教导少妇们善度贞洁的生活，爱丈夫，爱儿女，要明智、贞洁、勤理家务、善良、服从自己的丈夫，免得天主的话受到毁谤。}\par

有些缺点，老年有而青年没有。的确，有些是共有的，但此外老年人还有迟缓、胆怯、健忘、麻木和急躁。为此，保禄关于这些事劝勉老年人\Quote{要警醒。}因为在此时期有许多事使人不警醒，尤其是上文提到的，他们普遍麻木，难以激动。因此他又加上\Quote{庄重，审慎。}这里他指明智。因为审慎出自良好的心态。因为确实，有些老人狂热失态，或因酒，或因忧愁。因为老年使人气量狭窄。\par

\BibleQuote{在信德、爱德、忍耐上纯全无缺。}\par

他很好地加上了\Quote{在忍耐上}，因为这种美德特别适合老年人。\par

第3节。\BibleQuote{同样，劝老年妇人在行为上要合乎圣洁。} \BibleRef{铎 2:3}\par

就是说，她们在衣着和举止上要显出谦逊。\par

\BibleQuote{不诽谤人，不沉湎于酒。}\par

因为这是妇女和老年人的特有毛病。由于这个生命阶段天然的寒冷，产生了对酒的欲望，因此他将劝诫集中于此，以断绝一切醉酒的借口，愿她们远离此恶习，逃避随之而来的嘲笑。因为酒气更容易从下而上，脑膜因年老而受损，更易受伤害，这尤其导致醉酒。然而酒对这个年纪是必要的，因为其衰弱，但不需要很多。年轻妇女也不需要很多，虽然出于不同原因，因为酒点燃情欲之火。\par

\BibleQuote{是善事的导师。}\par

然而你禁止妇女教导，此处为何命令呢？因为你別处说：\BibleQuote{我不许女人施教} \BibleRef{弟前 2:12} 但请注意他加上的：\BibleQuote{也不许她管辖男人}。\par

第4节。\BibleQuote{好教导少妇们善度贞洁的生活} \BibleRef{铎 2:4}\par

请注意他如何将人们联系起来，如何使少妇服从年长的。因为他这里不是指女儿，而是仅指年龄。他意思是，每一位老年妇人要教导任何一位较年轻的妇人贞洁。\par

\BibleQuote{爱丈夫}\par

这是家庭一切美好的关键。\BibleQuote{男人和他的妻子和睦相处} \BibleRef{德 25:1}。因为哪里有这和睦，就没有不愉快的事。头与身体和谐，没有分歧，岂不其余肢体都平安？当统治者平安，谁还能分裂破坏和睦？另一方面，当他们彼此不和，家中就没有良好的秩序。因此这是最重要的事，比财富、地位、权力及其他一切都更重要。他不仅说平安，而且说\Quote{爱丈夫}。因为哪里有爱，就没有纷争，远远不会，反而会生出其他益处。\par

\BibleQuote{爱儿女}。这加得好，因为她若爱那根，就会更爱那果子。\par

\BibleQuote{要明智、贞洁、勤理家务、善良}。这一切都源于爱。由于对丈夫的爱，她们成为\Quote{勤理家务，善良}。\par

\BibleQuote{服从自己的丈夫，免得天主的话受到毁谤} \BibleRef{铎 2:5}\par

轻视丈夫的妇人，也疏忽自己的家；但出于爱则生出很大的节制，一切纷争都消失。如果丈夫是外教人，很快就会被说服；如果他是基督徒，会变得更善。你看到了保禄的屈尊吗？他凡事将我们引离世俗挂虑，这里却关心家务。因为当家务安排妥当时，就能为属灵的事留出空间，否则，属灵的事也会被损害。因为勤理家务的妇人也会节制，也会谨慎管理，不会倾向奢侈、不当花费等事。\par

\BibleQuote{免得天主的话受到毁谤} \BibleRef{铎 2:5}\par

你看他首先关心的是宣扬圣言，而非世俗之事；因为他写信给弟茂德时说：\BibleQuote{使我们可以敬虔庄重地过平静安宁的生活} \BibleRef{弟前 2:2}；而这里说\Quote{免得天主的话}和教义\Quote{受到毁谤}。因为如果一个信主的妇人与不信主的丈夫结婚，却不贤德，毁谤通常会波及天主；但如果她品行端正，福音便因她而得荣耀，也因她的善行。让那些与恶人或不信者结合的妇人听着，让她们学习以自己的榜样引导他们归向虔敬。因为即使你一无所获，未能吸引丈夫接受正确教义，至少你堵住了他的口，不容他毁谤基督教；这并非小事，而是大事，教义因我们的品行而受景仰。\par

第6节。\BibleQuote{同样劝年轻人要持重。} \BibleRef{铎 2:6}\par

请看他在如何处处推荐遵守\OriginalTerm{得体}{decorum}。因为他已将妇女教导的大部分托付给妇女，指派年长的教导年幼的。但男子的全部教导则交托给\Person{弟铎}本人。因为没有比克服非法享乐对这个年龄更困难的了。因为无论是爱财、求名，还是任何其他事物，都不如肉欲那样困扰年轻人。因此他略过其他事情，将其劝诫直指那个要害。但他并非要其他事情被忽略；因为他说什么呢？\par

第7节。\BibleQuote{在一切事上，显示自身为善行的模范。}\BibleRef{铎 2:7}\par

让年长妇女，他说，教导年轻的，但你自己劝勉年轻人要清醒自制。让你生活的光彩成为一所共有的学校、一个向公众展示的德性典范，好似一个原型模型，本身包含一切美善，提供榜样，使愿意的人可以轻易地将其中任何卓越印在自己身上。\par

第7节、第8节。\BibleQuote{在教导上显示不可败坏、庄重、真诚，以及无可指摘的正言，使那反对的一方羞愧，无可指责你们之处。}\par

他所说的\Quote{反对的一方}，是指魔鬼以及每个为他服务的人。因为当生命卓越，言语与之相应，温和良善，不给对手把柄时，将有难以言喻的益处。那么圣言的服务大有裨益，不是普通的言语，而是那被认可、无可指摘、不给愿意批判的人留下借口的言语。\par

第9节。\BibleQuote{劝勉仆人要服从自己的主人，凡事讨他们喜欢。}\BibleRef{铎 2:9}\par

你看到他先前所说的吗？\Quote{使那反对的一方羞愧，无可指责你们之处。}因此那借口\OriginalTerm{节欲}{continence}而分离妻子与丈夫的人，以及以其他藉口使仆人离开主人的人，都该受谴责。这不是\Quote{无可指摘的言语}，而是给不信者提供大把把柄，使众人反对我们。\par

\Quote{不可反驳。}\par

第10节。\BibleQuote{不可私取财物，反要显出全然的忠诚，好使他们凡事为天主我们的救主的教义增添光彩。}\BibleRef{铎 2:10}\par

因此他曾在另一处恰当地说：\Quote{做事如同事奉主，而非事奉人。}因为即使你善意地服侍你的主人，但这服侍的缘由来自你的敬畏，那以如此大的敬畏服侍祂的人，将获得更大的赏报。因为如果他不约束他的手或他桀骜不驯的舌头，外邦人怎能钦佩我们当中的教义呢？但如果他们看见他们的奴隶——那被教导基督哲学的人——比他们自己的哲学家更显自制，并以一切温良和善意服侍，他就会以各种方式钦佩福音的能力。因为希腊人不从教义本身判断教义，而是以生活和行为来检验教义。因此让妇女和仆人以他们的品行来做导师吧。因为不论在他们中间还是各处，都承认仆人的种族是激情洋溢的、不受感动、难以驾驭、不太容易接受德性教导，这不是由于他们的本性——愿天主不允许——而是由于他们没有教养，以及主人的疏忽。因为管治他们的人只关心自己的服侍；或者如果有时他们注意到仆人的道德，那也只是为了避免因仆人的私通、偷窃或酗酒给自己带来麻烦；仆人们就这样被忽视，无人关心他们，自然坠入邪恶的深渊。因为如果在一个父亲和母亲、一个监护人、一个主人和老师、合适的同伴、自由身份的荣耀以及其他许多优势的引导下，尚且难以避免与恶人亲近，那么对那些缺乏这一切、与恶人混在一起、随心所欲地交往而无人关心他们友谊的人，我们能期望什么呢？他们将是何等样的人？因此仆人很难成为良善的，尤其是当他们既没有从外人也没有从我们这里获得教导。他们不与有秩序的自由人交往，这些自由人大大注重自己的名声。综上所述，要有一个良善的仆人是一件困难且令人惊讶的事。\par

因此，当看到宗教的力量约束了本来自我意志强的阶层，使他们变得特别规矩和温和，他们的主人，无论多么不讲理，都会对我们的教义产生高度评价。因为很明显，他们先前在灵魂中灌输了复活、审判以及我们哲学教导死后所期待的一切恐惧，因而能够抵抗邪恶，在他们灵魂中有一个固定的原则来平衡罪的快感。所以\Person{保禄}如此关心这一类人并非偶然或无理由：因为他们越邪恶，悔改他们的宣讲力量就越令人钦佩。因为我们最钦佩医生的时候，是他将绝望的、无药可救的、无法控制荒谬欲望而沉溺其中的人恢复健康清醒。并且观察他最要求他们的是什么：那些最有助于他们主人安逸的品质。\par

\Quote{不可反驳，不可私取财物}；也就是说，在委托给它们的事上显出全部善意，特别忠信于主人的事务，并服从他们的命令。\par

德行。因此不要以为我无目的地详论此事。因为我以下的话将是对仆人说的。我亲爱的朋友，不要看你侍奉的是一个人，而要看你的侍奉是为了天主，是为福音增光。这样，你将怀着服从主人的心意去做一切，即使他急躁、无故发怒，你也容忍他。要想到你不是在讨好他，而是在履行天主的诫命；这样你就会轻易承受一切。我先前所说的，在此重复：当我们的灵性状态正确时，今生的事物也会随之而来。因为一个如此温顺、如此善良的仆人，不仅将蒙天主悦纳，分享那荣耀的冠冕，而且他的主人——即使他是粗野、铁石心肠、不人道、凶残的——也会称赞他、钦佩他，并尊崇他超过其余的人，让他管理众人，即使那主人是外邦人。\par

仆人须以这样的态度对待外邦主人，我将用一个例子向你们证明。若瑟与埃及人信仰不同，却被卖给了膳长。那么他做了什么呢？当他看见那年轻人品德高尚时，并不考虑信仰的差异，反而喜爱、恩待、钦佩他，并将其他人交给他管理，因为有了他，主人对自己家中的事一无所知。这样，若瑟成了第二个主人，甚至比他的主人更有权力，因为他比主人更了解主人的事务。后来，在我看来，当主人相信了妻子诬告的不实指控时，由于先前对若瑟的尊重，仍保持着对那义人的敬意，只以监禁来满足他的愤怒。因为如果主人不是因若瑟先前的品行而深深敬重他，就会用剑刺穿他的身体，立即将他处死。\BibleQuote{因为嫉恨是男人的愤怒；他决不宽容，也不因你多送礼物而罢休。} \BibleRef{箴 6:34}如果一般人的嫉恨尚且如此，那么作为埃及人和野蛮人，又自觉被他所尊荣的人伤害，他的愤怒岂不更加激烈？你们都知道，伤害并非来自所有人而对我们有同等影响，而是那些来自善待我们、信任我们且我们信任、并受过我们许多恩惠之人的伤害，最令我们痛苦和深刻。他没有如此思忖，也没有说：什么！我将一个仆人带进我的家，与他分享我的家产，使他自由，甚至使他比我更大，而他就是这样报答我的吗？他没有说这些话，因为他先前对他的尊重已占据了他的心。\par

若瑟在家中享有如此尊荣，这不足为奇；我们看到他在监狱中也得到了如此大的尊重。你们都知道监狱看守之人的心肠是何等惯于残忍。他们利用他人的不幸，折磨那些在患难中得到他人支持的人，用比野兽更残忍的方式从中牟利，这真是可悲。他们利用那些本应引起怜悯的悲惨境遇。我们还可注意到，他们并非以同样的方式对待所有囚犯；对于那些仅因控告而被监禁、且受冤屈的人，他们或许会怜悯；但对于那些因可耻和残暴罪行而被囚禁的人，他们则以无数酷刑惩罚。因此，典狱长不仅因这类人的习性而可能被期望是残忍的，也因若瑟被监禁的原因而更加残忍。因为谁不会对一个年轻人发怒呢？他蒙受如此大的尊荣，却被控以对主母图谋不轨来回报这些恩惠。典狱长想到这些，想到他所得的尊荣和所犯的罪过，难道不会以比野兽更残忍的方式对待他吗？但他因对天主的望德而超越这一切。因为灵魂的德行甚至能驯服野兽。若瑟以赢得主人欢心的同样温和，也征服了典狱长。这样，若瑟再次成为统治者：他在监狱中治理，如同在家中治理一样。因为他注定要统治，所以应当学习被治理；当他被治理时，他却成了治理者，并在家中作主。\par

因为如果保禄对提拔到教会的人要求说：\BibleQuote{人若不知道管理自己的家，怎能照管天主的教会呢？} \BibleRef{弟前 3:5}，那么那将要管理的人，首先应当是一个出色的管家。若瑟治理监狱，不像治理监狱，而像治理一个家。因为他减轻了所有人的灾难，并照顾那些被囚禁的人，如同照顾自己的肢体；不仅关心他们的不幸、安慰他们，而且若看见有人愁眉不展，他就上前询问原因，不能容忍看到任何人沮丧，或坐视不理，直到解除了他的沮丧为止。这种爱，许多人连对自己的子女也没有表现过。这些事可以追溯到他好运的开端。因为我们的努力必须先行，然后天主的祝福才会随之而来。\par

若瑟确实表现了这种关怀和关心，我们从故事中得知。据说，他看见法老的酒政和膳长被下在监里，就问：\BibleQuote{你们今日为什么面带愁容呢？} \BibleRef{创 40:7}不仅从这个问题，也从这些人的行为，我们可以看出他的美德。因为尽管他们是王的臣仆，却并不轻视他，也不在绝望中拒绝他的帮助，反而向他吐露全部秘密，如同向一位能同情他们的兄弟一样。\par

我说这一切，为证明：纵然有德者处于奴役、被掳、监禁，甚至地下深处，也没有任何事物能击败他。我对仆人们说这些，是要他们明白：即使他们有极其粗暴的主人——如同这位埃及人，或凶残如狱卒——他们仍能赢得主人的尊重；纵使主人是外邦人，不论如何，他们也能很快使之变得温良。因为没有什么比善行更吸引人，没有什么比温良、柔和与服从更令人愉悦和喜爱。具有这种品性的人能适应一切。他不以奴隶身份为耻，不回避穷人、病人和弱者。因为德行是卓越的，能胜过一切。如果德性在奴隶身上有如此力量，那么在自由人身上岂不更大？那么，让我们实践德行吧，无论为奴或自由，男人或女人。这样，我们必将蒙天主和人的喜爱；不仅为有德者所爱，也为恶人所爱；甚至更为恶人所爱，因为他们更尊敬并推崇德行。正如受管辖者最敬畏温良的人，邪恶者也最敬畏有德者，因为他们知道自己从何处坠落。既然德行的果实如此，让我们追求并得到它。如果我们持守德行，就没有什么可畏惧的，一切都将变得容易而可管理。即使我们历经水火，万物都为德行让路，甚至死亡本身。那么，让我们热切追求德行，好能获得未来的美善，在我们的主耶稣基督内，与祂一起，等等。\par

\SourceRecord{Translated by Philip Schaff.；Nicene and Post-Nicene Fathers, First Series；Vol. 13.；Edited by Philip Schaff.；Buffalo, NY: Christian Literature Publishing Co.,；1889.；<http://www.newadvent.org/fathers/23084.htm>.}

% structure: p0001 source=h1 target=section reason=page-title

\section{论 \BibleBook{弟铎}{书} 第五讲}

\BibleRef{铎 2:11-14}\par

\BibleQuote{因为天主的恩宠——带来救恩的——已显示给众人，教训我们，弃绝不虔敬和世俗的欲望，在今世清醒地、正义地、虔敬地生活，期待那有福的希望，以及伟大的天主和我们的救主耶稣基督光荣的显现。祂为我们舍弃了自己，为救赎我们脱离一切不义，洗净自己独特的子民，热心行善。}\par

既然从仆人们那里要求了如此伟大的德行，因为用一切行为装饰我们的天主和救主的道理是伟大的德行，并且嘱咐他们在普通事上也不要给主人任何冒犯的机会，他就添加了仆人们应当如此的正当事由：\BibleQuote{因为天主的恩宠——带来救恩的——已出现了。}那些以天主为师傅的人，很可能就是我所描述的那样，因为他们无数的罪过已获得赦免。因为你知道，除了其他考虑之外，这一点以不寻常的程度震慑并使灵魂谦卑：当它要为自己无数的罪过负责时，它没有受到惩罚，反而获得宽恕和无限的恩惠。因为如果有一个人，他的仆人犯了多起过错，他没有用鞭子抽打他，反而宽恕了他所有的过错，但要求他以后的行为要交账，并嘱咐他当心不要再犯同样的错误，而且赐予他极大的恩惠，你想，有谁听到这样的仁慈不会被打动呢？但不要以为恩宠只停留在赦免过去的罪过——它也确保我们将来不犯罪，因为这也是恩宠。因为如果祂从不惩罚那些仍然犯错的人，那就不是恩宠，而是鼓励恶行和邪恶。\par

\BibleQuote{因为天主的恩宠，}他说，\BibleQuote{出现了，教训我们，弃绝不虔敬和世俗的欲望，在今世清醒地、正义地、虔敬地生活，期待那有福的希望，以及伟大的天主和我们的救主耶稣基督光荣的显现。}你看，他是怎样把德行与报酬并列的。这也是恩宠，使我们脱离世俗之事，引领我们进入天堂。他在这里谈到两种显现；因为有两种：第一种是恩宠的显现，第二种是报应和公义的显现。\par

\BibleQuote{那弃绝不虔敬，}他说，\BibleQuote{和世俗的欲望。}\par

请看这里一切德行的根基。他没有说\Quote{避免，}而是说\Quote{弃绝。}弃绝意味着最远的距离、最大的憎恨和厌恶。他们以何等的决心和热忱转向脱离偶像，也要以同样的决心和热忱转向脱离恶习本身和世俗的欲望。因为这些也是偶像，即世俗的欲望和贪婪，他称之为偶像崇拜。一切对现世有用的事物都是世俗的欲望，一切与现世一同消亡的事物都是世俗的欲望。那么，我们不要与这些有任何关联。基督来了，\Quote{为使我们弃绝不虔敬。}不虔敬涉及教义，世俗的欲望涉及邪恶的生活。\par

\BibleQuote{并应在今世清醒地、正义地、虔敬地生活。}\par

你们看到，我一向肯定的，就是清醒不仅在于戒绝奸淫，而且我们必须摆脱其他的情欲。因此，爱财富的人并不清醒。因为如同奸淫者爱慕妇女，另一种人爱慕金钱，甚至更无节制，因为他并未被同样强烈的情欲所驱使。肯定地说，不能驾驭温驯马匹的车夫，比起不能约束野性难驯马匹的车夫，更是无能。那么怎样？他说，爱财富比爱妇女更弱吗？这从许多理由可以显明。首先，情欲来自本性的需要，而来自这种需要的东西必然难以约束，因为它植根于我们的本性。第二，因为古人并不看重财富，但他们对妇女在贞洁方面却非常看重。没有人指责按照法律与妻子同居到老的人，但所有人都指责囤积金钱的人。许多外邦哲学家轻视金钱，但没有一个人对妇女无动于衷，所以这种情欲比另一种更加专横。但是既然我们在向教会说话，让我们不要从外邦人而是从圣经中取例。因此，圣保禄几乎将此放在命令的等级中：\BibleQuote{只要有衣有食，就当知足。} \BibleRef{弟前 6:8} 但是关于妇女，他说：\BibleQuote{夫妻不可彼此亏负，除非两相情愿}——并\BibleQuote{再同房。} \BibleRef{格前 7:5} 你们常看到他制定关于合法结合的规则，他允许这种欲望的满足，允许再婚，对这事给予很多考虑，从未因之惩罚。但他到处指责贪爱金钱的人。关于财富，基督也常命令我们应避免它的腐蚀，但祂并没有说关于戒绝妻子的话。因为听祂关于金钱说的话：\BibleQuote{凡不撇下一切所有的，} \BibleRef{路 14:33} 但祂从未说：\BibleQuote{凡不撇下妻子的}，因为祂知道那情欲何等专横。圣保禄也说：\BibleQuote{婚姻，人人都当尊重，床也不可污秽} \BibleRef{希 13:4}；但他从未说过关心财富是值得尊重的，恰恰相反。因此他对弟茂德说：\BibleQuote{那些想要发财的人，就陷在迷惑、落在网罗和许多无知有害的私欲里} \BibleRef{弟前 6:9}。他没有说那些贪财的人，而是那些想要发财的人。\par

为使你从共同概念中学习这个问题的真实状态，必须从总体上向你呈现。如果一个男人一旦被剥夺了钱财，他就不再受其欲望的折磨，因为没有什么比拥有财富更能引起对财富的欲望了。但关于情欲并非如此，许多被阉割的人并没有摆脱内心燃烧的火焰，因为欲望存在于其他器官，深植于我们的本性之中。那么，为什么说这话呢？因为贪财者比行淫者更不节制，因为前者屈服于一种较弱的激情。确实，它更少来自激情，更多来自心灵的卑劣。但情欲是自然的，所以如果一个人不亲近女人，自然仍会履行其部分和运作。但在贪婪的情况下，却没有这类事。\par

\Quote{在今世度虔敬的生活。}\par

这希望是什么呢？我们劳苦的报酬是什么呢？\par

\Quote{期待所盼望的福分和显现。}\par

因为没有什么比那个显现更蒙福、更令人渴望了。言语无法描述它，其福乐远超我们的理解。\par

\Quote{期待所盼望的福分和我们伟大的天主及救主荣耀的显现。}\par

那些说圣子低于圣父的人在何处呢？\par

\Quote{我们伟大的天主及救主。}祂在我们还是仇敌时拯救了我们。那么，当我们已蒙悦纳时，祂还有什么不会做呢？\par

\Quote{伟大的天主。}当他说到天主是伟大时，他不是比较地说，而是绝对地说，因为在祂之后，没有谁是伟大的，因为这是相对的。因为如果是相对的，祂就是通过比较而伟大，而非本性的伟大。但现在，祂是无法比拟的伟大。\par

14节。\Quote{祂为我们舍弃了自己，为救赎我们脱离一切不法，并为自己炼净一个特别的子民。} \BibleRef{铎 2:14}\par

\Quote{特别的}：即从其余人中选出的，与他们毫无共通之处。\par

\Quote{热心行善。}\par

你看到我们的部分也是必要的吗？不仅是善行，而且是\Quote{热心}；我们应当以全部的喜悦和合宜的热忱在德行上向前迈进。因为当我们被罪恶重压、病入膏肓时，是祂的慈爱拯救了我们。但此后的事，既有祂的部分，也有我们的部分。\par

15节。\Quote{这些事你要讲论、劝勉，并用全部权柄责备。} \BibleRef{铎 2:15}\par

\Quote{这些事你要讲论、劝勉。}你看到他是如何命令弟茂德吗？\Quote{责斥、劝勉、责备。}但这里说：\Quote{用全部权柄责备。}因为这个民族的风俗更加顽固，所以他命令更严厉地责备他们，并用全部权柄。因为有些罪是应当用命令来防止的。我们可以用劝说劝人轻视财富、温良等。但奸淫者、行淫者、欺诈者，应当用命令使他们改过。那些沉迷于占卜和预言之类的人，应当\Quote{用全部权柄}加以纠正。请注意他是如何要他独立地、完全自由地坚持这些事的。\par

\Quote{不要让人轻看你。}但是\par

第三章1节。\Quote{你要提醒他们服从执政的和掌权的，听从官长，准备好行各样的善事，不诽谤人，不好争吵。} \BibleRef{铎 3:1}\par

那么，即使别人行恶，我们也不可辱骂他们吗？不只这样，还要\Quote{准备好行各样的善事，不诽谤人。}听这劝勉：\Quote{不诽谤人。}我们的嘴唇应当远离辱骂。因为即使我们的指责是真实的，也不应由我们来发出，而应由审判者来调查这事。因为他说：\Quote{你为什么判断你的兄弟呢？} \BibleRef{罗 14:10} 但若它们不真实，那火该有多大！听那强盗对他的同伴说的话：\Quote{我们也是同样的判决。} \BibleRef{路 23:40} 我们都在同样的危险中。你若辱骂别人，很快就会陷入同样的罪中。因此，蒙福的保禄劝告我们说：\Quote{那站着的，要小心，免得跌倒。} \BibleRef{格前 10:12}\par

\Quote{不好争吵，但要温和，向众人大显温柔。}\par

对希腊人和犹太人，对恶人和坏人。因为当他说：\Quote{那站着的，要小心，免得跌倒，}他是从未来唤醒他们的恐惧；但在这里，他相反地是从过去的考虑来劝勉他们，接下来的部分也是如此；\par

3节。\Quote{因为我们从前也是无知。} \BibleRef{铎 3:3}\par

他在致迦拉达人的书信中也这样做，在那里他说：\Quote{同样，我们从前为孩童的时候，受管于世界的基本元素。} \BibleRef{迦 4:4} 因此他说：不要诽谤任何人，因为你自己也曾如此。\par

\Quote{因为我们从前也是无知，悖逆，受迷惑，作各种私欲和享乐的奴隶，生活在恶毒和嫉妒中，是可憎的，又彼此仇恨。}\par

因此我们应当这样对待所有人，要有温柔的心。因为那从前处于这种状态、如今已得释放的人，不该指责别人，而该祈祷，感谢那施恩给他和他们、使他们脱离如此罪恶的主。没有人可以夸耀；因为众人都犯了罪。那么，若你行得好，却倾向于辱骂别人，当思想你自己从前的生活和未来的不确定，从而克制你的愤怒。因为即便你自幼年就度着有德的生活，你仍可能有许多罪；即便你认为自己没有，也要思想这不是你的德行的结果，而是天主的恩宠。因为若祂没有召叫你的祖先，你也会是悖逆的。请看这里他如何提到每一种邪恶。天主借着先知和所有其他方式行了多少事！我们可曾听见？\par

他说：\Quote{我们从前是受迷惑的。}\par

4节。\Quote{但当我们的救主天主的良善和慈爱显现时，}怎样呢？\Quote{不是由于我们自己所行的义行，而是照祂的怜悯，借着重生的洗礼和圣神的更新拯救了我们。} \BibleRef{铎 3:4}\par

真奇怪！我们是怎样沉溺于邪恶，以致不能得洁净，而需要重生呢？因为\Quote{重生}一词正暗示这一点。如同一所房子已濒临倒塌，没有人会在它下面撑支柱子，也不会在旧建筑上添补什么，而是将它拆到地基，重新建造。同样，天主并没有修补我们，而是将我们再造。这就是\Quote{圣神的更新}。祂使我们成为新人。如何？\Quote{藉着祂的圣神}；为进一步表明这点，他又加上，\par

第6节。\Quote{祂藉着我们的救主耶稣基督，丰厚地倾注在我们身上。}\BibleRef{铎 3:6}\par

因此，我们需要丰厚的圣神。\par

\Quote{好使我们因祂的恩宠成义}——再次强调是因恩宠而非因欠债——\Quote{按照永生的希望成为继承人。}\par

同时，这也激发谦逊和对未来的盼望。因为如果我们曾经如此被遗弃，以至于需要重生、因恩宠得救、心中毫无良善，而祂那时就拯救了我们，那么祂在未来的世界中更会拯救我们。\par

因为没有比基督来临前人类的残忍更糟糕的了。他们彼此相待如同敌人和战争。父亲杀死自己的儿子，母亲疯狂攻击儿女。没有秩序，没有自然法，也没有成文法；一切都颠倒了。常有奸淫、谋杀，甚至比谋杀更可怕的事，还有盗窃；事实上，一位外邦人告诉我们，这种行径被视为美德。这很自然，因为他们崇拜一位具有这种品性的神。他们的神谕经常要求他们杀死某些人。让我告诉你们当时的一个故事。米诺斯的儿子安德洛革俄斯来到雅典，在摔跤中获胜，因此受到惩罚并被处死。于是阿波罗以一种祸害补救另一种祸害，命令为此处死两倍七名青年。还有什么比这暴虐的命令更野蛮的？而且它被执行了。有人承担了平息那恶魔的疯狂，并杀死了这些青年，因为神谕的欺骗战胜了他们。但后来，当青年们抵抗并自卫时，此事便不再发生。如果当初是正义的，就不该被阻止；如果是不义的（它确实是不义的），就不该被命令。然后他们崇拜拳击手和摔跤手。他们不断进行战争，城与城、村与村、家与家。他们沉溺于男风，他们的一位智者制定法律，禁止奴隶从事少年爱和摔跤涂油，仿佛这是高贵的事；他们为此设有场所，公开进行。如果对他们所做的一切都加以叙述，就会看到他们公然违背自然，无人制止。他们的戏剧充斥着奸淫、淫荡和各种败坏。在他们不道德的夜间集会中，妇女被允许观看。可以看到一个处女在夜间坐在剧场中，周围是一群醉醺醺的青年疯狂狂欢。节日本身就是黑暗，以及他们实行的可憎行为。为此他说：\Quote{因为我们从前也曾经是愚昧的、悖逆的、受迷惑的、服事各样的私欲和宴乐。}一个人爱他的继母，一个女人爱她的继子，因此上吊自杀。至于他们对男童的激情，他们称之为\Quote{\OriginalTerm{少年爱}{Pædica}}，这种名字都不应提及。你们想看到儿子娶了母亲吗？这在他们中间也发生过，可怕的是，虽然出于无知，但他们崇拜的神并未阻止，反而允许这种对自然的凌辱发生，尽管她是一位贵妇。如果那些即使不为其他原因，也为了他们在众人中的名声而可能期望持守美德的人，都如此陷入恶行，那么大部分生活在默默无闻中的人会怎样呢？还有什么比这种快乐更多样化？某个人的妻子爱上了另一个男人，与奸夫合谋，在丈夫回来时杀死了他。你们大多数人可能知道这个故事。被杀者的儿子杀死了奸夫，然后杀了他的母亲，然后他自己疯了，被复仇女神纠缠。此后这个疯子又杀了另一个人，娶了他的妻子。有什么比这样的灾祸更糟糕？但我举这些外邦人的例子，是为了说服外邦人，当时世界上盛行了怎样的邪恶。但我们也可以从自己的经典中说明这一点。因为经上说：\Quote{他们将自己的儿女祭祀魔鬼。}\BibleRef{咏 105:37}另外，索多玛人除了因为违背自然的欲望外，没有其他原因而被毁灭。基督来临后不久，难道不是一位国王的女儿在醉酒的男人面前的宴会上跳舞，并要求以跳舞的奖赏杀死一位先知并取得他的头吗？\Quote{谁能述说上主的大能？}\BibleRef{咏 6:2}\par

他说：\Quote{可憎的，} \Quote{彼此憎恨。} 因为当我们放纵各种情欲在灵魂上时，必然会产生许多仇恨。因为哪里有与德行相连的爱，就没有人会欺骗别人。也要留意保禄所说的：\BibleQuote{你们不要自欺：无论是奸淫者、拜偶像者、通奸者、作娈童者、与男人行淫者、贪婪者、嗜酒者、辱骂者，都不能继承天主的国。你们中从前也有人是这样。} \BibleRef{格前 6:9} 你看，各种邪恶何等盛行？那是一片深沉的黑暗，一切正当都被败坏。因为那些享有预言恩惠，又亲眼看见许多灾祸降在敌人甚至自己身上的人，尚且不能自制，犯下无数愚昧的罪行，更何况其他人呢？他们的一个立法者下令，处女应在男人面前赤身摔跤。愿你们蒙受许多祝福！你们连听都受不了，但他们的哲学家却对实际行为不以为耻。另一位，他们的哲学首领，赞同他们出去打仗，并赞同女子成为公有，仿佛她是她们情欲的拉皮条者。\par

\Quote{生活在恶毒和嫉妒中。}\par

因为如果那些自称是哲学家的人制定了这样的法律，我们该怎么说那些不是哲学家的人呢？如果那些留着长胡子、披着庄重斗篷的人有这样的格言，其他人还有什么可说的？女人不是为此而被造的，人啊，决不是为了被当作公共娼妓。哦，你们这些颠覆一切礼仪的人，你们把男人当作女人来用，又把女人领出去作战，好像她们是男人！这是魔鬼的工作，要颠覆和混乱一切，越过从起初就设立的界限，并挪移天主为自然所定的界限。因为天主只将家务的照顾交给女人，而将公共事务的管理交给男人。但你们却将头降至脚，将脚抬到头。你们允许女人拿起武器，竟不以为耻。但我为何提这些呢？他们甚至在舞台上呈现一个杀害自己孩子的女人，而且不羞于用这样可憎的故事填满人的耳朵。\par

第4节。\BibleQuote{但当我们的救主天主的恩慈和爱显现时，祂救了我们，并不是由于我们自己所立的义行，而是依照祂的怜悯，藉着重生的洗礼和圣神的更新，这圣神是祂藉着我们的救主耶稣基督，丰丰富富地倾注在我们身上的，好使我们因祂的恩宠成义，而能按永生的望德成为继承人。} \BibleRef{铎 3:4}\par

\Quote{按照望德}是什么意思？是说，正如我们所希望的那样，我们将享受永生，或者因为你们甚至已经是继承人。\par

\Quote{这话是确实的。}\par

因为他一直谈论将来的事，而不是现在的事，所以他补充说，这话值得相信。他说，这些事是真实的，这一点从前面所说的就可显明。因为那位将我们从如此邪恶的状态和如此众多的邪恶中拯救出来的，也必定会将未来的美善赐予我们，如果我们居于恩宠中。因为一切都出自同样的慈爱关怀。\par

伦理训导：那么，让我们感谢天主，而不是辱骂他们；也不要控告他们，反而要恳求他们、为他们祈祷、劝告和忠告他们，即使他们侮辱和藐视我们。因为患病者的本性就是如此。但那些关心这些病人健康的人，会做一切事，忍受一切，即使可能没有效果，以免自己将来有疏忽之咎。难道你不知道吗？当医生对病人绝望时，旁边的一位亲属常常对他说：\Quote{再尽力诊治，不要留下任何未做的事，这样我就不必自责，就不会受责备，不会后悔。} 你难道没有看到近亲如何关心他们的亲属吗？他们为他们做多少事：恳求医生治愈他们，并坚持不懈地坐在他们旁边？让我们至少效法他们。然而，我们所关心的事物之间没有可比性。因为如果有人儿子身体患病，他绝不会拒绝长途跋涉来治愈他的疾病。但当灵魂处于糟糕状态时，没有人关心它；我们全都懒惰、漠不关心、疏忽大意，忽视我们的妻子、我们的孩子、以及我们自己，当我们被这种危险的疾病侵袭时。但等到为时已晚，我们才意识到。试想，事后说\Quote{我们从未预料到，我们从未想到会是这样}是多么可耻和荒谬。这不仅可耻，而且同样危险。因为在现世生活中，不为将来作准备是愚昧人的行为，更何况对于来世？我们听到许多人劝告我们，告诉我们什么该做，什么不该做。那么，让我们紧握那望德。让我们关心我们的救恩，让我们在一切事上呼求天主，愿祂向我们伸出援手。你还要懒惰多久？疏忽多久？我们对自己和同伴漠不关心要到几时？祂已丰丰富富地倾注了祂圣神的恩宠。因此，让我们思想祂赐予我们多么大的恩宠，并让我们自己也表现出同样大的热忱；或者，既然这不可能，也至少表现出一些，尽管较少。因为如果在接受这恩宠后我们仍麻木不仁，我们的惩罚将更重。祂说：\BibleQuote{因为如果我没有来向他们说话，他们就没有罪；但现在他们对自己的罪无可推诿。} \BibleRef{若 15:22} 但愿这样的话不会对我们说，并愿我们所有人都配得上那为爱祂之人所预备的福分，因我们的主耶稣基督，等等。\par

\SourceRecord{Translated by Philip Schaff.；Nicene and Post-Nicene Fathers, First Series；Vol. 13.；Edited by Philip Schaff.；Buffalo, NY: Christian Literature Publishing Co.,；1889.；<http://www.newadvent.org/fathers/23085.htm>.}

% structure: p0001 source=h1 target=section reason=page-title

\section{论《弟铎书》第六篇讲道}

\BibleRef{铎 3:8-11}\par

\Quote{我要你切实讲论这些事，好让那些已信从天主的人，用心行善。这些事是美好的，且对人有益。但要远避愚拙的辩论、\OriginalTerm{族谱}{genealogies}、\OriginalTerm{争竞}{contentions}以及关于法律的争执，因为它们是无益且虚妄的。一个\OriginalTerm{异端者}{heretic}在第一次和第二次劝诫之后，就要弃绝他。因为你知道，这样的人已经背道，并且犯罪，自己定了自己的罪。}\par

在论到天主对人的爱、祂对我们说不尽的眷顾、我们从前是什么以及祂为我们所成就的事之后，他接着就说：\Quote{我要你切实讲论这些事，好让那些已信从天主的人，用心行善}\BibleRef{铎 3:8}，也就是说，讲论这些事，并借着对这些事的默想来劝人施舍。因为前面所说的话不仅适用于谦卑、不自高自大、不辱骂别人，也适用于其他一切德行。同样，在与格林多人辩论时，他说：\Quote{你们知道我们的主本是富有的，却成了贫穷的，好使我们因祂的贫穷而成为富有的。}\BibleRef{格后 8:9}在考虑了天主对人的眷顾和极深的爱之后，他便由此劝勉他们施舍，而且不是以平常和轻微的方式，而是\Quote{他们\Emph{用心}行善}，也就是说，不仅要资助受害者，不仅用金钱，还要用庇护和保护，要护卫寡妇和孤儿，为所有受苦的人提供避难所。这便是行善。他说：\Quote{这些事是美好的，且对人有益。但要远避愚拙的辩论、\OriginalTerm{族谱}{genealogies}、\OriginalTerm{争竞}{contentions}以及关于法律的争执，因为它们是无益且虚妄的。}\BibleRef{铎 3:9}这些\OriginalTerm{族谱}{genealogies}是什么意思？因为在他写给弟茂德的信中也提到\Quote{无稽的传说和无穷的族谱}\BibleRef{弟前 1:4}。［或许这里和那里都是指犹太人，他们以亚巴郎为祖先而自夸，却忽略了自己的本分。因此他称它们为\Quote{愚拙}和\Quote{无益}；因为信赖无益的事是愚拙的。］\Quote{争竞}，他指的是与异端者的争论，他不想我们徒劳无益地劳苦，因为从中一无所获，并且毫无结果。因为当一个人已经背道，并且下定决心无论如何都不改变心意时，你何必白费力气在石头上播种，而应当把你的尊贵辛劳用在自己的百姓身上，与他们谈论施舍和其他一切德行呢？那么，他在别处怎么说\Quote{也许天主会赐给他们悔改}\BibleRef{弟后 2:25}，但在这里却说\Quote{一个异端者在第一次和第二次劝诫之后，就要弃绝他，因为你知道，这样的人已经背道，并且犯罪，自己定了自己的罪}\BibleRef{铎 3:10-11}呢？在前一处，他是指着那些他抱有希望的、仅仅提出反对的人说的。但当他已被众人所知道和显明时，你又何必枉自争辩呢？何必击打空气呢？\Quote{自己定了自己的罪}是什么意思？因为他不能说没有人告诉过他，没有人劝诫过他；因此，既然劝诫之后他依然如故，他就是自己定自己的罪了。\par

第12节。\Quote{我打发\Person{阿尔特玛}或\Person{提希苛}到你那里去的时候，你要赶紧到我这里来，到\Place{尼科波利}去。}\BibleRef{铎 3:12}你说什么？你既已派他管理克里特，怎么又召他到你那里去呢？这并不是要把他从工作中抽离，而是要更好地训练他。因为他并不是要他来伺候自己，好像把他到处带在身边作为随从一样，这从他下面的话可以看出来：\par

\Quote{因为我决定在那里过冬。}\par

\Place{尼科波利}是\Place{色雷斯}的一个城市。\par

第14节。\Quote{你要赶紧送\Person{则纳}律师和\Person{阿颇罗}上路，使他们一无所缺。}\BibleRef{铎 3:14}\par

这些人不是受托管理教会的人，而是他的同伴。但\Person{阿颇罗}更为激昂，\Quote{他是有口才的人，又精通圣经}\BibleRef{宗 18:24}。你也许会说，如果\Person{则纳}是律师，就不应由别人来供养。但这里的\Quote{律师}是指通晓犹太法律的人。他似乎说：丰丰富富地供应他们的需要，使他们什么都不缺。\par

第14至15节。\Quote{我们的人也要学习行善，以备必需之用，免得成为不结果实的。同我在一起的众人都问候你。请代我问候那些在\OriginalTerm{信德}{faith}中爱我们的人。}\par

也就是说，要么是那些爱\Person{保禄}本人的人，要么是那些忠信的人。\par

\Quote{愿恩宠与你们众人同在。阿们。}\par

那么，你怎能命令他堵住反对者的口，如果他在那些人自取灭亡时必须放过他们呢？他的意思是，他不可主要为了他们的益处而做，因为他们既已心志败坏，就不会从中获益。但如果他们伤害了别人，他就应该起来与他们争辩，勇敢地等候他们；但如果你被逼无奈，看见他们毁灭别人，就不要沉默，而要堵住他们的口，是出于对将被他们毁灭之人的关心。一个热心而生活正直的人确实不可能避免争论，但要像我说的那样去做。因为这种恶来自闲散和虚妄的哲学，即只关心空谈。当一个人应当教导、祈祷或感谢时，却说出过多的话语，这是极大的损害。因为吝惜我们的钱财而不吝惜我们的话语是不对的；我们更应该节约话语，甚于节约钱财，而不应当将我们自己交给各色各样的人。\par

\Quote{要他们谨慎于维持善工}何意？即不要等待贫乏者前来，而要主动寻找需要帮助的人。如此，谨慎之人显出他的关切，并会以极大的热忱履行此责。因为在行善时，接受恩惠者所获之益，不及行善者所得之利，因这恩赐使他们在天主前获得信心。但在另一情况下，争论不止；因此他称异端者不可救药。因为若忽略那些尚有悔改希望的人属懈怠，则对病入膏肓者徒费心力实为极愚极狂；因这反使他们更加胆大。\par

他说：\Quote{愿我们的人也学习为必需之用途而维持善工，免得他们不结果实。}你注意到他为我们自己的人比那些接受他们恩惠的人更为焦虑。因为他们或许会由许多其他人护送前行，但他说：我关切的是我们自己的朋友。因为他人若掘出宝藏并供养他们的导师，这对他们有何益处？他们仍不结果实。那么，难道基督不能用五饼喂饱五千人，用七饼喂饱四千人，祂不能供养自己和门徒吗？\par

伦理训导。那么，祂为何由妇女供养？因为经上说：\BibleQuote{妇女跟随祂，并服事祂。}\BibleRef{谷 15:41}这是要我们从起初就学会：祂关切那些行善的人。难道保禄自己亲手供养他人，就不能不靠他人帮助而维持生计吗？但你见他既接受又请求援助。且听他如此行的理由：\BibleQuote{我并不求馈赠，而是求归入你们账上的丰硕果实。}\BibleRef{斐 4:17}起初，当人们变卖所有财产放在宗徒脚前时，你看见宗徒们更为关心他们，而非那些接受施舍的人。因为若他们只关心穷人无论如何得到救济，就不会如此严厉审判阿纳尼雅和撒斐辣扣留钱财的罪。保禄也不会吩咐人\BibleQuote{不要忧愁，也不要出于不得已}\BibleRef{格后 9:7}而施与。你说什么，保禄？你劝阻施与穷人吗？不，他回答：我不仅考虑他们的益处，也考虑施与者的好处。你看见，当先知给拿步高那卓越的忠告时，他并非只考虑穷人。因他不满足于只说\Quote{施与穷人}，而是说什么？\BibleQuote{以施舍赎你的罪，以怜恤穷人赎你的不义。}\BibleRef{达 4:27}舍弃你的财富，并非为使他人得饱，而是为叫你免于惩罚。基督又说：\BibleQuote{去变卖你所有的，分给穷人……然后来跟随我。}\BibleRef{玛 19:21}你看见这命令是为使他能跟随祂？因财富是阻碍，故此祂命他将财富分给穷人，以教导灵魂慈悲怜悯，轻视财富，远避贪吝。因学会施与需要者的人，终将学会不接受当施与者的给予。这使人肖似天主。然而，守贞、禁食、卧地苦修比这更困难，但没有什么比施舍更能熄灭我们罪恶之火。它超越所有其他德行。它将爱施舍者置于君王本身之侧，实为公义。因为守贞、禁食、卧地苦修的功效仅限实行者本人，无人由此得救。但施舍广及众人，怀抱基督的肢体；那些影响多人的行为远大于仅限于一人者。\par

因为施舍是爱之母，这爱是基督徒的特征，超越一切奇迹，基督的门徒由此彰显。它是我们罪恶的良药，灵魂污秽的洗涤，直通天堂的阶梯；它连结基督的身体。你愿学习此事何等卓越？在宗徒时代，人们变卖财产带到他们面前，然后分配。因为经上说：\BibleQuote{照每人所需的分给各人。}\BibleRef{宗 4:35}请告诉我，且不谈论将来，也不考虑那将来的天国，我们且看现世生活中，谁更有收获：是接受者，还是施与者？前者抱怨并彼此争吵。后者却一心一意。经上说：\BibleQuote{他们一心一意，并且众人都蒙受了恩宠。}\BibleRef{宗 4:32}他们在极大的纯朴中生活。你看见他们因施与而获益了吗？请告诉我，你愿与谁同列：是与那些给出财产而一无所有的人，还是与那些甚至接受他人财物的人？\par

看施舍的果实：分隔和障碍被移除，他们的灵魂立刻结合在一起。\Quote{他们全都一心一意，灵魂合一。}因此，即使撇开施舍不谈，舍弃财富也能带来益处。我说这些话，是叫那些没有从祖先继承遗产的人不致沮丧，好像他们比富人更为贫穷。因为他们若愿意，他们拥有的更多。因为他们会更乐意施舍，像那寡妇一样，并且他们不会对邻人有任何仇视，而且各方面都享有自由。这样的人不会被没收财产的威胁，他凌驾于一切不义之上。正如那些不被衣服累赘而奔跑的人不易被抓住，而那些被许多衣服和长袍拖累的人很快就被追上；同样，富人和穷人也是如此。一个即使被抓住，也会轻易逃脱；而另一个即使不被拘留，也被自己的绳索——无数的忧虑、痛苦、情欲、激怒——拖累，这些全都压制灵魂。不仅如此，还有财富所带来许多其他事物。富人比穷人更难节制和俭朴生活，更难摆脱情欲。那么，你又说，他将获得更大的赏报。——绝非如此。——什么？难道他克服了更大的困难，不就更当得吗？——但这些困难是他自找的。因为我们没有被命令致富，反而是相反。但他为自己准备了如此多的绊脚石和障碍。\par

另有些人不仅舍弃财富，还刻苦身体，如同行窄路的旅客。你非但不这样做，反而更炽热地燃烧你情欲的火炉，并且使你身边聚集更多。所以，去吧，进入宽路，因为那正是接待像你这样之人的路。但窄路是为那些受苦困迫的人，他们随身携带的只有那些能穿过窄路的负担，如施舍、仁爱、良善和温良。如果你背负这些，你将容易找到入口；但如果你携带骄傲、被情欲燃烧的灵魂，以及那荆棘的负担——财富，那么你需要宽大的空间才能通过，而且你很难不撞到他人并在路上压倒他们而进入人群。在这种情况下，你需要与他人保持很远的距离。但携带金银的人——我指的是德行的成就——不会使邻人逃离他，反而使人更靠近他，甚至与他结合。但如果财富本身是荆棘，那么贪婪又该是什么呢？你为什么要把那东西带走？是为了给那火添柴，使火焰更大吗？地狱之火还不够吗？想想三少年是如何战胜火窑的。把那火窑想象为地狱。他们在患难中被投入其中，捆绑着锁链；但在里面他们找到了宽裕之地；而周围站着的人则不然。\par

类似的事现在也要经历，如果我们勇敢地抵抗围绕我们的考验。如果我们对天主怀有希望，我们就会安然无恙，拥有宽广的余地，而那些使我们陷入困境的人必将灭亡。因为经上记载：\BibleQuote{谁掘洞穴，必陷其中。}\BibleRef{德 26:27}尽管他们捆绑我们的手脚，但苦难将有力量使我们解脱。因为请看这个奇迹：那些被人捆绑的人，火却释放了他们。好比有些人被交到他们朋友的仆人手中，仆人因尊重主人的友谊，非但不加害，反而十分尊重他们；同样，火知道三少年是自己上主的朋友，就挣断了他们的锁链，释放了他们，让他们离去，并且成为他们脚下的铺路石，被他们践踏。这很合理，因为他们是为天主的荣耀而被投入其中的。我们这些受苦难的人，务要持守这些榜样。\par

但是，你看，他们说他们从患难中被拯救了，而我们却没有。确实，他们得到了拯救，而且合理，因为他们进入那火窑时并不期望得救，而像是准备去死。请听他们所说的话：\BibleQuote{天上有一位天主，他必拯救我们。但若不然，大王，你当知道，我们决不侍奉你的神，也不朝拜你所立的金像。}\BibleRef{达 3:17}但我们好像在与上主的惩戒讨价还价，甚至定下时间，说：\Quote{若他到这时还不施怜悯……}因此我们没有得救。\Person{亚巴郎}离开家时，绝不是期望再次得到儿子，而是准备祭献他；他出乎意料地重新收到了安全的儿子。你也是，当你陷入患难时，不要急于得救，要准备忍受一切，这样你很快就会从苦难中得救。因为天主使你受苦正是为此，即要惩戒你。因此，当我们从一开始就学会耐心承受，而不陷入绝望时，他立刻就会解救我们，因为整个事情已经成就。\par

我愿给你们讲一个有益的故事，其中含有很多益处。是什么呢？有一次，当一场迫害兴起，一场激烈的战争正攻击教会时，有两个人被逮捕。一个人已经准备好忍受任何事；另一个人虽坚定愿意被斩首，但恐惧战栗地退缩不敢忍受其他酷刑。请观察天主对待这些人的安排。当法官就坐时，他下令将那个愿意忍受一切的人斩首。另一个人则被命令吊起来受刑，而且不止一次或两次，而是从一个城市到另一个城市。为什么允许这样呢？为使他通过痛苦恢复他所忽略的那种心态，使他摆脱一切的怯懦，不再害怕忍受任何事情。同样，\Person{若瑟}急切地想逃离监狱时，却被留在那里。因为请听他所说的话：\Quote{我确是从希伯来地被偷来的；但请你向国王提起我。}\BibleRef{创 40:14}为此他被允许留在那里，为使他学会不将希望或信赖寄托于人，而将一切交托于天主。知道了这些事，让我们感谢天主，并让我们做一切有益于我们的事，为使我们获得将来的美善，藉着我们的主耶稣基督，愿光荣归于圣父，偕同圣神，从今直到永远，世世无尽。阿们。\par

\SourceRecord{Translated by Philip Schaff.；Nicene and Post-Nicene Fathers, First Series；Vol. 13.；Edited by Philip Schaff.；Buffalo, NY: Christian Literature Publishing Co.,；1889.；<http://www.newadvent.org/fathers/23086.htm>.}
